\documentclass[11pt]{article}

\usepackage[letterpaper,margin=0.82in,headheight=15pt]{geometry}
\usepackage[T1]{fontenc}
\usepackage{lmodern}
\usepackage{microtype}
\usepackage{amsmath,amssymb,amsthm,bm,mathtools}
\usepackage{booktabs,longtable,tabularx,array,multirow}
\usepackage{enumitem}
\usepackage{xcolor}
\usepackage{fancyhdr}
\usepackage{graphicx}
\usepackage{tikz-cd}
\usepackage{hyperref}
\usepackage[nameinlink,noabbrev]{cleveref}

\definecolor{navy}{HTML}{17324D}
\definecolor{teal}{HTML}{147D92}
\definecolor{orange}{HTML}{D47A27}
\definecolor{softblue}{HTML}{EEF4F7}
\definecolor{softgray}{HTML}{F3F5F6}
\definecolor{darkgray}{HTML}{46545E}
\definecolor{softorange}{HTML}{FFF4E8}

\hypersetup{
  colorlinks=true,
  linkcolor=navy,
  citecolor=teal,
  urlcolor=teal,
  pdftitle={Volume I: Regulated Light-Front State Spaces and Microscopic Dynamics},
  pdfauthor={DeuteronWigner project}
}

\pagestyle{fancy}
\fancyhf{}
\fancyhead[L]{\small\color{darkgray}DeuteronWigner formalism - Volume I}
\fancyhead[R]{\small\color{darkgray}Regulated light-front foundations}
\fancyfoot[L]{\small\color{darkgray}28 July 2026}
\fancyfoot[R]{\small\color{darkgray}\thepage}
\renewcommand{\headrulewidth}{0.4pt}

\setlist[itemize]{leftmargin=1.5em,itemsep=2pt,topsep=3pt}
\setlist[enumerate]{leftmargin=1.7em,itemsep=2pt,topsep=3pt}
\setlength{\parindent}{0pt}
\setlength{\parskip}{5pt}
\setlength{\emergencystretch}{2em}
\renewcommand{\arraystretch}{1.16}
\allowdisplaybreaks

\newtheorem{definition}{Definition}[section]
\newtheorem{axiom}[definition]{Axiom}
\newtheorem{principle}[definition]{Design principle}
\newtheorem{requirement}[definition]{Requirement}
\newtheorem{proposition}[definition]{Proposition}
\newtheorem{theorem}[definition]{Theorem}
\newtheorem{lemma}[definition]{Lemma}
\newtheorem{corollary}[definition]{Corollary}
\newtheorem{remark}[definition]{Remark}

\crefname{definition}{definition}{definitions}
\crefname{axiom}{axiom}{axioms}
\crefname{principle}{principle}{principles}
\crefname{requirement}{requirement}{requirements}
\crefname{proposition}{proposition}{propositions}
\crefname{theorem}{theorem}{theorems}
\crefname{lemma}{lemma}{lemmas}
\crefname{corollary}{corollary}{corollaries}
\crefname{remark}{remark}{remarks}

\newcommand{\kT}{\bm{k}_{T}}
\newcommand{\pT}{\bm{p}_{T}}
\newcommand{\qT}{\bm{q}_{T}}
\newcommand{\PT}{\bm{P}_{T}}
\newcommand{\DT}{\bm{\Delta}_{T}}
\newcommand{\DNT}{\bm{\Delta}_{N T}}
\newcommand{\bD}{\bm{b}_{\Delta}}
\newcommand{\bM}{\bm{b}_{\mathrm{TMD}}}
\newcommand{\RT}{\bm{R}_{T}}
\newcommand{\dd}{\mathrm{d}}
\newcommand{\Tr}{\operatorname{Tr}}
\newcommand{\tr}{\operatorname{tr}}
\newcommand{\Span}{\operatorname{span}}
\newcommand{\End}{\operatorname{End}}
\newcommand{\Hom}{\operatorname{Hom}}
\newcommand{\Obj}{\operatorname{Obj}}
\newcommand{\Id}{\operatorname{Id}}
\newcommand{\Ker}{\operatorname{Ker}}
\newcommand{\Ran}{\operatorname{Ran}}
\newcommand{\Spec}{\operatorname{Spec}}
\newcommand{\diag}{\operatorname{diag}}
\newcommand{\Hil}{\mathcal H}
\newcommand{\Kin}{\mathcal K}
\newcommand{\Phys}{\mathcal H_{\mathrm{phys}}}
\newcommand{\Alg}{\mathcal A}
\newcommand{\Rep}{\mathcal V}
\newcommand{\Path}{\mathsf{Path}}
\newcommand{\LF}{\mathrm{LF}}
\newcommand{\TMD}{\mathrm{TMD}}
\newcommand{\GTMD}{\mathrm{GTMD}}
\newcommand{\GPD}{\mathrm{GPD}}
\newcommand{\CP}{\mathrm{CP}}
\newcommand{\Amp}{\mathbf{Amp}}
\newcommand{\Dens}{\mathbf{Dens}}
\newcommand{\Match}{\mathbf{Match}}
\newcommand{\Red}{\mathbf{Red}}
\newcommand{\Proc}{\mathbf{Proc}}
\newcommand{\cO}{\mathcal O}
\newcommand{\cM}{\mathcal M}
\newcommand{\cR}{\mathcal R}
\newcommand{\cS}{\mathcal S}
\newcommand{\cV}{\mathcal V}
\newcommand{\cW}{\mathcal W}
\newcommand{\cE}{\mathcal E}
\newcommand{\cG}{\mathcal G}
\newcommand{\cQ}{\mathcal Q}
\newcommand{\cP}{\mathcal P}
\newcommand{\cC}{\mathcal C}
\newcommand{\cZ}{\mathcal Z}
\newcommand{\cD}{\mathcal D}
\newcommand{\cN}{\mathcal N}
\newcommand{\cF}{\mathcal F}
\newcommand{\cA}{\mathcal A}
\newcommand{\cB}{\mathcal B}
\newcommand{\eps}{\varepsilon}
\newcommand{\vark}{\varkappa}
\newcommand{\bpsi}{\bm\psi}
\newcommand{\bbone}{\bm 1}
\newcommand{\comm}[2]{[#1,#2]}
\newcommand{\acom}[2]{\{#1,#2\}}
\newcommand{\abs}[1]{\left|#1\right|}
\newcommand{\norm}[1]{\left\lVert#1\right\rVert}
\newcommand{\inner}[2]{\left\langle #1\middle|#2\right\rangle}

\newcolumntype{Y}{>{\raggedright\arraybackslash}X}
\newcolumntype{L}[1]{>{\raggedright\arraybackslash}p{#1}}

\title{\vspace{-1.0cm}
{\color{navy}\bfseries Volume I: Regulated Light-Front State Spaces\\[2mm]
and Microscopic Dynamics}\\[4mm]
\large Hilbert-space construction, gauge constraints, Fock-sector transitions,\\
renormalized mass operators, and controlled truncation}
\author{DeuteronWigner project}
\date{28 July 2026}

\begin{document}
\maketitle

\begin{center}
\begin{minipage}{0.94\textwidth}
\colorbox{softblue}{\parbox{0.96\textwidth}{
\textbf{Status and purpose.}
This is the first dynamical volume following the architecture specification
in Volume~0.  It is a normative formalism: it defines the regulated state
spaces, physical constraints, mass operator, renormalization datum,
truncation system, current consistency conditions, and typed computational
interfaces that must exist before a light-front eigenstate can serve as the
common microscopic parent of quark, antiquark, and gluon GTMDs.  It does not
select a final numerical Hamiltonian or claim that a finite truncation is
continuum QCD.  The Stage~0 repository audit at commit
\texttt{5d4641f} is treated as the immutable regression baseline for the
adapter migration described in \cref{sec:c1-spine}.
}}
\end{minipage}
\end{center}

\begin{abstract}
A predictive spin-1 partonic model requires more than a complete registry of
TMD projections.  It requires a regulated light-front quantum state from
which flavor, helicity, orbital interference, sea, gluon, current, and later
nuclear observables descend through shared operators.  This volume gives the
foundational formalism for that state.  We fix the front-form spacetime and
normalization conventions; distinguish intrinsic momentum, transfer, Wigner,
and TMD Fourier variables as separate mathematical types; construct
regulated one-particle and Fock-sector spaces with exact particle statistics;
define the gauge-constrained physical Hilbert space and color-singlet
subspaces; separate exact gauge constraints, approximate internal symmetries,
and the dynamical restoration of rotations; and formulate the regulated
mass-squared operator as a self-adjoint block operator containing canonical,
effective, counterterm, and sector-induced interactions.

A complete model datum includes the regulator, basis, retained Fock sectors,
zero-mode and boundary prescription, counterterms, renormalization
conditions, currents, and operator scheme.  Eliminating sectors is shown to
induce both effective interactions and effective operators.  Calculations at
different cutoffs form a directed comparison system rather than necessarily
a literal chain of nested Hilbert spaces; convergence is defined for matched
observables, while bare Fock probabilities remain resolution dependent.  We
give normalization, charge, momentum, helicity, and angular-momentum ledgers;
Ward-identity and spin-1 angular-condition diagnostics; eigensolver and state
tracking requirements; and a machine-oriented type spine for coordinates,
transverse rank and mass normalization, sectors, Wilson paths, decorated
operators, and the distinct map classes \(\Amp,\Dens,\Match,\Red,\Proc\).
The final sections state propositions, property-based tests, acceptance
criteria, and the exact contract that a state must satisfy before entering
Volume~II, where common nonzero-transfer quark and gluon GTMD operators will
be constructed.
\end{abstract}

\tableofcontents

\section{Role, scope, and formal status}
\label{sec:role}

\subsection{Normative formalism versus repository audit}

Volume~0 specified the architecture of the next model class.  The Stage~0
Codex task then audited the existing repository without changing its accepted
physics.  This volume has a different logical status:
\begin{equation}
\boxed{\text{Volume I is normative}}
\qquad\text{whereas}\qquad
\boxed{\text{the Stage 0 audit is descriptive and diagnostic}.}
\label{eq:normative-descriptive}
\end{equation}
The definitions below are therefore not inferred from legacy class names.
Instead, the legacy model is wrapped by explicit adapters until its data can
be migrated without changing any authoritative artifact.

\begin{requirement}[Regression immutability]
\label{req:regression-immutability}
Until a later work package explicitly authorizes a physics change, the eight
authoritative parent/correlator artifacts and every accepted numerical value
at commit \texttt{5d4641f} are immutable.  A type migration is accepted only
if it is byte-identical on those artifacts and passes the complete baseline
validation suite.
\end{requirement}

\subsection{Stage 0 baseline incorporated into this volume}

The Stage~0 report supplied for this volume records the following baseline.
The counts supersede earlier documentation counts where the repository gained
additional tests during the audit.

\begin{longtable}{L{0.29\textwidth}L{0.61\textwidth}}
\toprule
Item & Stage~0 result at commit \texttt{5d4641f}\\
\midrule
\endhead
Repository regression & 484 of 484 tests passed.\\
Acceptance builders & Every documented acceptance builder passed.\\
Evidence ledger & 36 of 36 evidence rows passed.\\
Atlas rendering & 162 of 162 atlas pages rendered.\\
Numerical identity & All eight authoritative parent/correlator hashes were unchanged.\\
Main architecture gaps & Fragmented operator identity; generic transverse-coordinate aliases; untyped composition across \(\Amp,\Dens,\Match,\Red,\Proc\); incomplete gluon path and \(f/d\) color identity; no enforceable provenance exclusion graph.\\
Convention conclusion & The implemented light-front convention is internally consistent and no conflicting accepted tensor sign was found.  The principal risks are implicit specialized Fourier formulas and missing transverse-rank and mass-normalization metadata.\\
\bottomrule
\end{longtable}

These findings determine the implementation emphasis of \cref{sec:c1-spine};
they do not alter the physical definitions developed in the intervening
sections.

\subsection{Relation to the canonical boundary and later volumes}

The accepted canonical model is a complete declared leading-twist forward
boundary at its stated scale and grid.  Its central state is a retained
parton-target correlator assembled from fitted and modeled scalar amplitudes,
realistic deuteron structure, gauge-link scenarios, and typed uncertainty
axes.  It remains the comparison and regression oracle.  The next model
changes the direction of inference:
\begin{equation}
\{\cM_{
m ren}^2,\theta_{\rm micro},\cR\}
\longrightarrow
\{|\Psi_N\rangle,|\Psi_D\rangle,U_\gamma\}
\longrightarrow
\{W_{a/N},W_{a/D}\}
\longrightarrow
\{F_A,\cO_{\rm process}\}.
\label{eq:generative-chain-volume1}
\end{equation}
Volume~I defines the first arrow up to the construction of regulated states
and consistent local operators.  The remaining volumes are separated as
follows:
\begin{enumerate}
\item \textbf{Volume II:} common quark, antiquark, and gluon GTMD operators,
nonzero-transfer overlaps, and commuting TMD/GPD/PDF/current reductions;
\item \textbf{Volume III:} Wilson-line dynamics, absorptive phases, and
naive-\(T\)-odd color/link sectors;
\item \textbf{Volume IV:} matched nucleon-to-deuteron light-front theory and
spin-1 nuclear amplitudes;
\item \textbf{Volume V:} regulator-to-QCD matching, rank-aware evolution,
and process factorization;
\item \textbf{Volume VI:} common-parameter inference, correlated
uncertainties, and withheld validation.
\end{enumerate}

\subsection{Status vocabulary}

\begin{definition}[Formal status labels]
Every assertion or computational result shall be classified as one of:
\begin{description}[leftmargin=2.45cm,style=nextline]
\item[Exact] A consequence of the declared algebra, representation,
kinematics, operator definition, or finite-dimensional identity.
\item[Regulator-exact] Exact inside a specified regulated theory or basis,
but not asserted to survive removal of the regulator without matching.
\item[Scheme-defined] Fixed by a declared renormalization, rapidity,
operator, or normalization scheme.
\item[Controlled] The leading term of an approximation with a declared
expansion parameter and convergence or remainder test.
\item[Calibrated] Determined from a restricted set of data through shared
microscopic parameters.
\item[Model] Physically motivated but underdetermined, with explicit
replacement and sensitivity semantics.
\item[Numerical] A finite precision, quadrature, discretization, eigensolver,
or interpolation statement with a measured tolerance.
\end{description}
\end{definition}

Passing an exact or numerical test does not upgrade a model assumption to a
first-principles result.  Conversely, an approximation does not become
uncontrolled merely because it is effective, provided its domain, matching,
and remainder are explicit.

\subsection{Nonclaims}

This volume does not claim:
\begin{itemize}
\item a solved continuum QCD spectrum;
\item that one finite Hamiltonian matrix is regulator independent;
\item that global color-singlet projection alone enforces every local gauge
constraint;
\item that light-front Fock probabilities are scheme-independent
observables;
\item that a valence-only nucleon can predict antiquark or gluon GTMDs;
\item that a tensor network supplies interactions omitted from the mass
operator;
\item that an effective nuclear Hamiltonian may be combined with unrelated
partonic operators without matching;
\item that a successful type migration changes the evidential status of the
canonical boundary.
\end{itemize}

\section{Front-form spacetime, normalization, and intrinsic variables}
\label{sec:front-form}

\subsection{Coordinates, metric, and distributions}

We use
\begin{equation}
 x^{\pm}=\frac{x^0\pm x^3}{\sqrt2},
 \qquad
 p\cdot x=p^+x^-+p^-x^+-\bm p_T\cdot\bm x_T,
\label{eq:lf-coordinates}
\end{equation}
with metric \(g^{+-}=g^{-+}=1\), \(g^{ij}=-\delta^{ij}\), and
\(\epsilon^{0123}=+1\).  Hence
\begin{equation}
 p^2=2p^+p^- -\bm p_T^2.
\label{eq:lf-invariant}
\end{equation}
The light-front hypersurface is \(x^+=0\).  We write
\begin{equation}
 \delta^{(3)}_{\LF}(p-p')
 =\delta(p^+-p^{\prime +})\,
  \delta^{(2)}(\bm p_T-\bm p'_T),
\qquad
 [\dd p]=\frac{\dd p^+\,\dd^2\bm p_T}
 {2p^+(2\pi)^3}.
\label{eq:lf-measure}
\end{equation}
Only positive-\(p^+\) particle modes are included in the ordinary Fock
basis.  Any longitudinal zero mode is a separate constrained or dynamical
object governed by the zero-mode datum of \cref{sec:regulator-datum}; it is
not silently approximated by the endpoint of a positive-\(p^+\) quadrature.

\begin{requirement}[Convention closure]
\label{req:convention-closure}
Every state, kernel, transform, current, and observable must declare the
light-front convention of \cref{eq:lf-coordinates,eq:lf-invariant} or pass
through an explicit adapter.  No function may infer a factor of \(2\),
\(\sqrt2\), or \(2\pi\) from a variable name.
\end{requirement}

\subsection{One-particle normalization}

For a species label \(a\), internal label \(\alpha\), and helicity
\(\lambda\), the continuum one-particle convention is
\begin{equation}
 \langle p',a',\alpha',\lambda'|p,a,\alpha,\lambda\rangle
 =2p^+(2\pi)^3\delta^{(3)}_{\LF}(p-p')
 \delta_{aa'}\delta_{\alpha\alpha'}\delta_{\lambda\lambda'}.
\label{eq:one-particle-normalization}
\end{equation}
Creation and annihilation operators obey
\begin{align}
 \acom{b_{q\alpha\lambda}(p),b^\dagger_{q'\alpha'\lambda'}(p')}
 &=2p^+(2\pi)^3\delta^{(3)}_{\LF}(p-p')
 \delta_{qq'}\delta_{\alpha\alpha'}\delta_{\lambda\lambda'},
\label{eq:quark-algebra}\\
 \acom{d_{q\alpha\lambda}(p),d^\dagger_{q'\alpha'\lambda'}(p')}
 &=2p^+(2\pi)^3\delta^{(3)}_{\LF}(p-p')
 \delta_{qq'}\delta_{\alpha\alpha'}\delta_{\lambda\lambda'},
\label{eq:antiquark-algebra}\\
 \comm{a_{c\lambda}(p)}{a^\dagger_{c'\lambda'}(p')}
 &=2p^+(2\pi)^3\delta^{(3)}_{\LF}(p-p')
 \delta_{cc'}\delta_{\lambda\lambda'}.
\label{eq:gluon-algebra}
\end{align}
All other elementary (anti)commutators vanish.  A discretized basis replaces
Dirac deltas by Kronecker deltas and includes the quadrature weights in the
basis isometry, not in later physics kernels.

\subsection{Poincare generators in the front form}

The ten generators are \(P^\mu\) and \(M^{\mu\nu}=-M^{\nu\mu}\).  On the
surface \(x^+=0\), the seven kinematical generators are
\begin{equation}
 P^+,\quad P^1,\quad P^2,\quad M^{12},\quad M^{+-},
 \quad M^{+1},\quad M^{+2},
\label{eq:kinematical-generators}
\end{equation}
and the dynamical generators are
\begin{equation}
 P^-,\qquad M^{-1},\qquad M^{-2}.
\label{eq:dynamical-generators}
\end{equation}
The interacting invariant mass-squared operator is
\begin{equation}
 \cM^2
 =P_\mu P^\mu
 =2P^+P^- -\PT^2.
\label{eq:mass-operator}
\end{equation}
In a frame with \(\PT=0\), solving \(P^-\) at fixed \(P^+\) is equivalent
to solving \(\cM^2\), but the project stores \(\cM^2\) explicitly so that
the factor of two and transverse subtraction cannot drift.

The Pauli--Lubanski vector is
\begin{equation}
 W^\mu=-\frac12\epsilon^{\mu\nu\rho\sigma}M_{\nu\rho}P_\sigma,
\label{eq:pauli-lubanski}
\end{equation}
and light-front helicity is defined on a massive irreducible state by
\begin{equation}
 \Lambda=\frac{W^+}{P^+}.
\label{eq:lf-helicity}
\end{equation}
At finite basis or Fock truncation, \(J^z=M^{12}\) can be exact while the
transverse rotation generators remain approximate.  The formalism therefore
does not infer full rotational covariance from an exact \(J^z\) block label.

\subsection{Intrinsic constituent variables}

For an \(N\)-particle sector of total momentum \(P\), define
\begin{equation}
 p_i^+=x_iP^+,
 \qquad
 \bm p_{iT}=x_i\PT+\bm k_{iT},
\label{eq:intrinsic-momenta}
\end{equation}
with
\begin{equation}
 x_i>0,
 \qquad
 \sum_{i=1}^N x_i=1,
 \qquad
 \sum_{i=1}^N\bm k_{iT}=0.
\label{eq:intrinsic-constraints}
\end{equation}
The intrinsic momentum manifold is
\begin{equation}
 \cM_N=
 \left\{(x_i,\bm k_{iT})_{i=1}^N
 \;\middle|\;\cref{eq:intrinsic-constraints}\right\}.
\label{eq:intrinsic-manifold}
\end{equation}
A convenient project measure is
\begin{equation}
 [\dd\Gamma_N]
 =\left[\prod_{i=1}^N
 \frac{\dd x_i\,\dd^2\bm k_{iT}}{16\pi^3}\right]
 16\pi^3\,
 \delta\!\left(1-\sum_i x_i\right)
 \delta^{(2)}\!\left(\sum_i\bm k_{iT}\right).
\label{eq:intrinsic-measure}
\end{equation}
The factors \(x_i^{-1/2}\) associated with the canonical relativistic
normalization are placed in the Fock expansion in
\cref{eq:fock-expansion}, not duplicated in \cref{eq:intrinsic-measure}.
Other standard conventions are admissible only through a declared unitary or
isometric normalization adapter.

\begin{proposition}[Kinematical boost invariance of intrinsic variables]
\label{prop:intrinsic-boost}
Under a common kinematical longitudinal or transverse light-front boost, the
variables \(x_i\) and \(\bm k_{iT}\) of
\cref{eq:intrinsic-momenta} are invariant.
\end{proposition}

\begin{proof}
A longitudinal boost rescales every \(p_i^+\) and \(P^+\) by the same
positive factor, leaving \(x_i=p_i^+/P^+\) invariant.  A transverse
kinematical boost shifts \(\bm p_{iT}\mapsto\bm p_{iT}+x_i\bm\beta_T P^+\)
and \(\PT\mapsto\PT+\bm\beta_TP^+\).  Their difference
\(\bm p_{iT}-x_i\PT=\bm k_{iT}\) is unchanged.
\end{proof}

\subsection{Typed transverse variables}
\label{sec:typed-transverse}

Equal dimensions do not imply equal physical meaning.  The project uses the
following distinct transverse types.

\begin{longtable}{L{0.14\textwidth}L{0.22\textwidth}L{0.28\textwidth}L{0.25\textwidth}}
\toprule
Symbol & Software kind & Physical meaning & Fourier or recoil relation\\
\midrule
\endhead
\(\DT\) & \texttt{DELTA\_T} & External transverse momentum transfer to a hadron or target. & Conjugate to \(\bD\).\\
\(\bD\) & \texttt{B\_DELTA} & Partonic impact coordinate in the transverse Wigner transform. & \(\bD\leftrightarrow\DT\).\\
\(\kT\) & \texttt{K\_T} & Average intrinsic transverse momentum of the active parton. & Conjugate to the bilocal separation \(\bM\).\\
\(\bM\) & \texttt{B\_TMD} & TMD impact/separation variable used in matching and evolution. & \(\bM\leftrightarrow\kT\).\\
\(\DNT\) & \texttt{DELTA\_NT} & Transfer assigned to an active nucleon in a nuclear overlap. & Related to \(\DT\) by an explicit recoil map.\\
\(\pT\) & \texttt{P\_T\_NUCLEAR} & Intrinsic transverse momentum of a constituent nucleon in a nuclear state. & Conjugate to a nuclear relative coordinate only after a declared transform.\\
\(\RT\) & \texttt{R\_T\_NUCLEAR} & Nuclear transverse Wigner coordinate. & \(\RT\leftrightarrow\DNT\).\\
\(\qT\) & \texttt{Q\_T\_MEASURED} & Measured transverse imbalance in a process. & Appears in process Fourier kernels and is not a parton-state coordinate.\\
\bottomrule
\end{longtable}

\begin{axiom}[No transverse aliasing]
\label{ax:no-transverse-alias}
The types in the preceding table are pairwise noninterchangeable.  An
operation accepting one type may receive another only through a named map
whose formula, frame, Jacobian, sign, and provenance are part of its identity.
\end{axiom}

\subsection{Fourier transforms and transverse rank metadata}
\label{sec:rank-metadata}

For the conventions used later,
\begin{align}
 \rho_W(x,\kT,\bD)
 &=\int\frac{\dd^2\DT}{(2\pi)^2}
 e^{-i\bD\cdot\DT}W(x,\kT,\DT),
\label{eq:wigner-fourier}\\
 \widetilde F(x,\bM)
 &=\int\dd^2\kT\,e^{+i\bM\cdot\kT}F(x,\kT).
\label{eq:tmd-fourier}
\end{align}
A scalar radial transform is not sufficient for nonzero transverse rank.  If
an irreducible harmonic has angular weight \(m\), define a rank convention
\begin{equation}
 \mathfrak r=
 \bigl(m,\tau_m,N_m(k_T;M_{\rm ref}),
 \widetilde N_m(b_{\rm TMD};M_{\rm ref}),\chi_m\bigr),
\label{eq:rank-spec}
\end{equation}
where \(\tau_m\) names the tensor basis, \(N_m\) and \(\widetilde N_m\)
record extracted momentum and mass powers, and \(\chi_m\) records the Fourier
phase.  The radial pair is then represented schematically by
\begin{align}
 \widetilde F^{(m)}(b)
 &=2\pi\chi_m
 \int_0^\infty k\,\dd k\,J_{|m|}(bk)
 N_m(k;M_{\rm ref})F^{(m)}(k),
\label{eq:rank-forward-transform}\\
 F^{(m)}(k)
 &=\frac{\chi_m^{-1}}{2\pi}
 \int_0^\infty b\,\dd b\,J_{|m|}(bk)
 \widetilde N_m(b;M_{\rm ref})\widetilde F^{(m)}(b).
\label{eq:rank-inverse-transform}
\end{align}
The functions \(N_m\) and \(\widetilde N_m\) are not inferred from \(m\);
they depend on whether powers such as \(k_T/M\), \(bM\), and \(i^m\) have
been factored into the named scalar coefficient or retained in the tensor.

\begin{requirement}[Rank-complete identity]
\label{req:rank-complete}
Every transverse coefficient must carry \(m\), tensor-basis identity,
reference mass, extracted powers, Fourier sign, Fourier phase, and Bessel
order.  A rank-aware adapter may reproduce a legacy specialized formula; a
generic scalar transform may not guess it.
\end{requirement}

\section{Regulator, basis, and zero-mode datum}
\label{sec:regulator-datum}

\subsection{The regulator label}

A regulated calculation is indexed by a structured label
\begin{equation}
 r=\bigl(
 \cB, K,N_{\max},L_{\max},\Lambda_{\rm UV},\Lambda_{\rm IR},
 \cF_r,\epsilon_{\rm grid},\epsilon_{\rm quad},
 \cG_r^{\rm gauge},\cZ_r^{(0)},\mathfrak p_r
 \bigr).
\label{eq:regulator-label}
\end{equation}
Here \(\cB\) is the one-particle basis family; \(K\) denotes longitudinal
resolution when a discretized longitudinal basis is used; \(N_{\max}\) and
\(L_{\max}\) denote transverse or angular cutoffs when applicable;
\(\Lambda_{\rm UV,IR}\) are physical or derived regulator scales;
\(\cF_r\) is the retained Fock-sector set; \(\epsilon_{\rm grid,quad}\) are
numerical resolution labels; \(\cG_r^{\rm gauge}\) specifies gauge fixing and
residual constraints; \(\cZ_r^{(0)}\) specifies the zero-mode prescription;
and \(\mathfrak p_r\) contains boundary conditions and any additional
regulator metadata.

\begin{definition}[Complete regulated theory datum]
\label{def:regulated-theory-datum}
At level \(r\), the theory datum is
\begin{equation}
 \mathfrak T_r=
 \bigl(
 \Hil_r,\cD_r,\cM_r^2,
 \{\cO_{A,r}\},
 \cR_r,\bm c_r,\bm\cC_r,
 \mathfrak s_r,\mathfrak g_r
 \bigr),
\label{eq:regulated-theory-datum}
\end{equation}
where \(\cD_r\) is the common invariant operator domain,
\(\cR_r\) the regulator prescription, \(\bm c_r\) the bare/effective
coefficient vector, \(\bm\cC_r\) the renormalization conditions,
\(\mathfrak s_r\) the renormalization and operator scheme, and
\(\mathfrak g_r\) the gauge/zero-mode/boundary datum.
\end{definition}

A matrix file without the remaining fields of
\cref{eq:regulated-theory-datum} is a numerical object, not a defined
Hamiltonian model.

\subsection{One-particle species spaces}

Let the species set be
\begin{equation}
 \mathsf A=\{q_f,\bar q_f,g\mid f\in\mathsf F_{\rm active}\}.
\label{eq:species-set}
\end{equation}
For each \(a\in\mathsf A\), define the regulated one-particle kinematical
space
\begin{equation}
 \Hil_a^{(r)}
 =\cB_a^{(r)}
 \otimes V_a^{\rm flavor}
 \otimes V_a^{\rm color}
 \otimes V_a^{\rm hel}.
\label{eq:one-particle-space}
\end{equation}
The momentum/basis factor \(\cB_a^{(r)}\) is an \(L^2\) space or a finite
quadrature/basis approximation with a declared isometry into the continuum.
The internal factors are
\begin{align}
 V_{q_f}^{\rm color}&\cong\bm 3,
 &V_{\bar q_f}^{\rm color}&\cong\overline{\bm 3},
 &V_g^{\rm color}&\cong\bm 8,
\label{eq:color-reps}\\
 V_{q_f}^{\rm hel}&=\Span\{|+\tfrac12\rangle,|-\tfrac12\rangle\},
 &V_g^{\rm hel}&=\Span\{|+1\rangle,|-1\rangle\}
\label{eq:helicity-reps}
\end{align}
for physical transverse gluon modes after the constrained components have
been treated.  If regulator fields, longitudinal polarizations, ghosts, or
auxiliary species are retained, they inhabit a separately tagged enlarged
kinematical space and are never silently identified with physical partons.

Orbital angular momentum is generated on \(\cB_a^{(r)}\) or on the
multi-particle intrinsic momentum space.  It is not an independent physical
tensor factor unless the chosen basis explicitly diagonalizes an orbital
operator and stores the corresponding basis label.

\subsection{Continuum, basis, and quadrature realizations}

A basis realization supplies maps
\begin{equation}
 \iota_r:\Hil_r^{\rm basis}\longrightarrow\Hil^{\rm cont},
 \qquad
 \pi_r:\Hil^{\rm cont}\longrightarrow\Hil_r^{\rm basis},
\label{eq:basis-isometries}
\end{equation}
with \(\pi_r\iota_r=\Id_r\) for an orthonormal finite basis, or a measured
quadrature defect for a collocation/grid representation.  A basis coefficient
array has meaning only together with:
\begin{itemize}
\item basis-family and normalization identifiers;
\item ordering of every mode and internal label;
\item quadrature weights or orthonormality metric;
\item ultraviolet and infrared scales implied by the truncation;
\item phase conventions for helicity and angular basis states;
\item a serialization version and content hash.
\end{itemize}

\begin{requirement}[No coefficient-only state]
\label{req:no-coefficient-only}
No coefficient array may be interpreted as a state without its basis and
inner-product metric.  Dot products, normalizations, and adjoints must be
computed with the declared metric.
\end{requirement}

\subsection{Longitudinal zero modes and constrained fields}

The project shall not assume that the ordinary positive-\(p^+\) Fock basis
contains every degree of freedom relevant to symmetry breaking, gauge
constraints, or boundary Wilson lines.  The zero-mode datum
\(\cZ_r^{(0)}\) must state one of the following, with supporting equations:
\begin{enumerate}
\item zero modes are retained as explicit constrained or dynamical variables;
\item they are integrated out and their induced operators are included;
\item boundary conditions prove their irrelevance for the declared
observable and truncation;
\item their omission is an explicit model discrepancy with a validation test.
\end{enumerate}
A bare statement that the light-front vacuum is trivial is not a zero-mode
prescription.

\section{Fock sectors, statistics, and the graded Hilbert space}
\label{sec:fock-spaces}

\subsection{Canonical occupation basis}

For flavor-resolved occupation numbers, define a Fock label
\begin{equation}
 \nu=\bigl(
 \{n_{q_f}\}_f,\{n_{\bar q_f}\}_f,n_g;
 B,Q,S,\ldots;J^z,\Pi_{\LF},\chi
 \bigr),
\label{eq:fock-label}
\end{equation}
where the first entries specify particle content and the remaining entries
specify exact or controlled block quantum numbers.  Derived charges must
agree with the occupations.  The label \(\chi\) is reserved for additional
sector identities such as regulator species or nuclear-resolution sectors.

The preferred computational representation is a canonical occupation-number
basis.  Fermionic antisymmetry and bosonic symmetry are then enforced by the
operator algebra.  If a first-quantized tensor-product representation is
used, define the statistics projector
\begin{equation}
 \cP_{\nu}^{\rm stat}
 =\prod_{s\in\mathrm{fermions}}\cA_s
  \prod_{s\in\mathrm{bosons}}\cS_s,
\label{eq:statistics-projector}
\end{equation}
where \(\cA_s\) and \(\cS_s\) antisymmetrize and symmetrize identical species,
respectively.

\begin{requirement}[Factorials occur once]
\label{req:factorials-once}
Identical-particle factorials shall occur either in the normalized occupation
basis or in the first-quantized statistics projector and integration measure,
but not in both.  The serialization manifest records which convention is
used.
\end{requirement}

\subsection{Intrinsic Fock expansion}

With the one-particle normalization of
\cref{eq:one-particle-normalization}, a hadron state is expanded as
\begin{equation}
\begin{split}
 |P,\Lambda;h\rangle_r
 ={}&\sum_{\nu\in\cF_r}
 \sum_{\bm\alpha_\nu}
 \int[\dd\Gamma_{N_\nu}]
 \frac{\psi^{(r),\Lambda}_{h,\nu}
 (\{x_i,\bm k_{iT},\alpha_i\})}
 {\sqrt{\prod_{i=1}^{N_\nu}x_i}}
 \\
 &\times
 |\nu;\{x_iP^+,x_i\PT+\bm k_{iT},\alpha_i\}\rangle,
\end{split}
\label{eq:fock-expansion}
\end{equation}
where \(\bm\alpha_\nu\) collects flavor, color, helicity, and discrete basis
labels and the state ket uses the normalized occupation convention of the
project.  The wave functions satisfy
\begin{equation}
 \sum_{\nu\in\cF_r}
 \sum_{\bm\alpha_\nu}
 \int[\dd\Gamma_{N_\nu}]
 \left|\psi^{(r),\Lambda}_{h,\nu}
 (\{x_i,\bm k_{iT},\alpha_i\})\right|^2=1
\label{eq:fock-normalization}
\end{equation}
for a unit-normalized internal state.  The exact placement of
\(x_i^{-1/2}\) factors is convention dependent, but once
\cref{eq:fock-expansion} is chosen it must be used in every vertex, current,
overlap, and sum rule.

\subsection{Kinematical and physical sector spaces}

The kinematical sector is
\begin{equation}
 \Kin_\nu^{(r)}
 =\cP_\nu^{\rm stat}
 \left[
 \bigotimes_{i=1}^{N_\nu}\Hil_{a_i}^{(r)}
 \right]_{\rm intrinsic},
\label{eq:kinematical-sector}
\end{equation}
where the subscript enforces
\cref{eq:intrinsic-constraints}.  Gauge and global-quantum-number constraints
select the physical sector
\begin{equation}
 \Hil_{\nu,\rm phys}^{(r)}
 =\cP_{\nu}^{\rm phys,(r)}\Kin_\nu^{(r)}.
\label{eq:physical-sector}
\end{equation}
The complete regulated physical space is
\begin{equation}
 \Phys^{(r)}
 =\bigoplus_{\nu\in\cF_r}\Hil_{\nu,\rm phys}^{(r)}.
\label{eq:graded-physical-space}
\end{equation}
Let \(P_\nu^{(r)}\) denote the orthogonal sector projector.

\begin{axiom}[Orthogonal sector grading]
\label{ax:orthogonal-sectors}
For the retained sector set,
\begin{equation}
 P_\nu P_\mu=\delta_{\nu\mu}P_\nu,
 \qquad
 P_\nu^\dagger=P_\nu,
 \qquad
 \sum_{\nu\in\cF_r}P_\nu=\Id_r.
\label{eq:sector-projector-algebra}
\end{equation}
\end{axiom}

\begin{proposition}[Norm ledger]
\label{prop:norm-ledger}
For a normalized state \(|\Psi\rangle\in\Phys^{(r)}\), the sector weights
\begin{equation}
 Z_\nu^{(r)}=\langle\Psi|P_\nu^{(r)}|\Psi\rangle
\label{eq:sector-weight}
\end{equation}
are nonnegative and satisfy \(\sum_\nu Z_\nu^{(r)}=1\).
\end{proposition}

\begin{proof}
Orthogonal projectors are positive, so \(Z_\nu\ge0\).  Using completeness in
\cref{eq:sector-projector-algebra},
\(\sum_\nu Z_\nu=\langle\Psi|\Id_r|\Psi\rangle=1\).
\end{proof}

\begin{remark}[Resolution dependence]
The identity in \cref{prop:norm-ledger} is regulator-exact.  Individual
\(Z_\nu^{(r)}\) generally change with the regulator, basis, unitary
transformation, and sector definition.  They are diagnostics of the chosen
resolution, not automatically observables.
\end{remark}

\subsection{Minimal nucleon sector content}

The first nucleon model intended to generate valence, sea, and gluon
correlators must include at least
\begin{equation}
 |N,\Lambda\rangle
 =\psi_{qqq}^{\Lambda}|qqq\rangle
 +\psi_{qqqg}^{\Lambda}|qqqg\rangle
 +\psi_{qqqq\bar q}^{\Lambda}|qqqq\bar q\rangle
 +\cdots.
\label{eq:minimal-nucleon-fock}
\end{equation}
The ellipsis is not a license to hide missing physics.  Every omitted sector
must be associated with either an induced operator, a power-counted remainder,
or an explicit truncation uncertainty.  The retained basis must permit the
helicity and orbital interferences required by the later transverse-rank
inventory, including \(L_z=0,\pm1,\pm2\) where the named projections demand
them.

\begin{requirement}[No unsupported correlator sector]
\label{req:no-unsupported-sector}
A central antiquark or gluon correlator may not be generated from a state that
contains no corresponding dynamical or matched degree of freedom.  A
comparison input may be attached externally, but it remains a comparison
object and may not be labeled a microscopic prediction of the state.
\end{requirement}

\section{Gauge constraints, color, and the physical subspace}
\label{sec:gauge}

\subsection{Gauge fixing is not the physical-state condition}

The kinematical color tensor product is larger than the physical hadron
space.  A gauge choice such as \(A^+=0\) removes redundancy from the
coordinate description, but it does not by itself replace Gauss's law,
residual gauge constraints, boundary data, or color-singlet conditions.  At
\(x^+=0\), define the regulated Gauss operator schematically by
\begin{equation}
 \cG_r^a(x)
 =\bigl(D_\mu F^{\mu +}\bigr)^a_r(x)
 -g_r\,\bar\psi_r(x)\gamma^+t^a\psi_r(x)
 +\cG_{\rm reg,r}^a(x),
\label{eq:gauss-operator}
\end{equation}
where \(\cG_{\rm reg,r}^a\) contains regulator, auxiliary-field, or boundary
contributions required by the chosen formulation.

\begin{definition}[Gauge-consistency datum]
\label{def:gauge-datum}
The gauge datum \(\mathfrak g_r\) comprises:
\begin{enumerate}
\item the gauge condition and treatment of constrained field components;
\item the regulated Gauss operators or BRST charge;
\item residual gauge transformations and boundary conditions at light-front
infinity;
\item the zero-mode prescription;
\item the physical-state condition;
\item the Wilson-line boundary convention used by later nonlocal operators;
\item Ward, Slavnov--Taylor, or BRST identities used as finite-truncation
diagnostics.
\end{enumerate}
\end{definition}

In a direct physical-state construction, the physical subspace obeys
\begin{equation}
 \cG_r^a[f]|\Psi\rangle=0
 \quad\text{for all admissible test functions }f,
\label{eq:gauss-physical-state}
\end{equation}
up to a measured regulator defect.  In a BRST formulation it is replaced by
cohomology of a nilpotent regulated BRST charge.  The implementation must
choose one route; it may not combine the state space of one route with the
inner product of another without an explicit equivalence map.

\subsection{Constrained light-front fields}

In light-front gauge, nondynamical spinor and gauge components are solved by
constraint equations.  Substituting those solutions into the canonical
Hamiltonian produces instantaneous interactions.  Accordingly, the
interaction inventory must include, where generated by the chosen gauge and
regularization:
\begin{equation}
 \cV_{\rm can}
 =\cV_{qqg}+\cV_{3g}+\cV_{4g}
 +\cV_{\rm inst,f}+\cV_{\rm inst,g}
 +\cV_{\rm mass}.
\label{eq:canonical-interaction-classes}
\end{equation}
The labels in \cref{eq:canonical-interaction-classes} denote operator
classes, not fitted TMD components.  Omitting an instantaneous term because
it is inconvenient is a Hamiltonian truncation and must appear in the error
and Ward-identity ledgers.

\subsection{Global color and singlet projection}

Let the total regulated color generator be
\begin{equation}
 Q_r^a
 =\sum_{i\in q}t_i^a
 -\sum_{j\in\bar q}(t_j^a)^*
 +\sum_{k\in g}(T_{\rm adj}^a)_k
 +Q_{\rm reg}^a.
\label{eq:total-color-generator}
\end{equation}
A physical isolated hadron satisfies
\begin{equation}
 Q_r^a|\Psi_h\rangle=0
\label{eq:color-singlet-condition}
\end{equation}
within the regulator tolerance.  Denote the singlet projector in sector
\(\nu\) by \(\cP_{\nu}^{\bm1}\).  Examples include the baryon tensor
\(\epsilon_{ijk}/\sqrt6\), color traces in \(q\bar q\) and \(gg\) channels,
and nontrivial singlet recouplings in multi-quark sectors.

\begin{proposition}[Invariant operators preserve singlets]
\label{prop:singlet-preservation}
If an operator \(O_r\) commutes with every total color generator,
\(\comm{Q_r^a}{O_r}=0\), then
\begin{equation}
 O_r\Ran\cP^{\bm1}\subseteq\Ran\cP^{\bm1}.
\label{eq:singlet-preservation}
\end{equation}
\end{proposition}

\begin{proof}
For a singlet \(|\Psi\rangle\),
\(Q_r^aO_r|\Psi\rangle=O_rQ_r^a|\Psi\rangle=0\).  Thus
\(O_r|\Psi\rangle\) lies in the simultaneous kernel of the color generators,
which is the singlet subspace under the declared finite representation.
\end{proof}

\begin{remark}[Singlet is necessary, not sufficient]
Global color-singlet projection does not prove local gauge consistency.  The
Gauss, residual-gauge, zero-mode, and boundary requirements in
\cref{def:gauge-datum} remain independent acceptance conditions.
\end{remark}

\subsection{Color intertwiners and vertex identity}

A color coupling is an intertwiner, not an anonymous numerical tensor.  For
representations \(R_{\rm in}\) and \(R_{\rm out}\),
\begin{equation}
 C\in\Hom_{SU(3)}(R_{\rm in},R_{\rm out})
 \quad\Longleftrightarrow\quad
 C D_{\rm in}(g)=D_{\rm out}(g)C.
\label{eq:color-intertwiner}
\end{equation}
The canonical quark-gluon, three-gluon, and four-gluon vertices use the
corresponding \(t^a\), \(f^{abc}\), and products of structure constants with
all normalization conventions declared.  For later gluon GTMD and TMD
operators, the link pair and the antisymmetric \(f\)-type versus symmetric
\(d\)-type color class are part of operator identity.  They may not be
reconstructed from a process name after the correlator has been evaluated.

\begin{requirement}[Color identity completeness]
\label{req:color-identity}
Every state tensor and operator tensor carrying color indices must declare:
its representation on each index, incoming/outgoing variance, invariant
coupler, normalization, basis ordering, and singlet or open-color status.
Every gluon nonlocal operator must additionally declare the ordered link pair
and any \(f/d\) color-contraction class.
\end{requirement}

\subsection{Finite-truncation gauge diagnostics}

At each regulator level, report separately:
\begin{align}
 \epsilon_{\rm singlet}^{(r)}
 &=\max_a\norm{Q_r^a|\Psi\rangle},
\label{eq:singlet-defect}\\
 \epsilon_{\rm Gauss}^{(r)}
 &=\sup_{a,f\in\mathfrak F_r}
 \frac{\norm{\cG_r^a[f]|\Psi\rangle}}
 {\norm{|\Psi\rangle}},
\label{eq:gauss-defect}\\
 \epsilon_{\rm Ward}^{(r)}
 &=\sup_{q\in\cQ_r}
 \frac{\abs{q_\mu\langle P'|J_r^\mu|P\rangle}}
 {\cN_J(P',P,q)},
\label{eq:ward-defect-preview}
\end{align}
with declared test sets \(\mathfrak F_r,\cQ_r\) and normalization
\(\cN_J\).  A small global singlet defect cannot be substituted for a
missing Ward or Gauss test.

\section{Spin, orbital structure, and representation grammar}
\label{sec:symmetry}

\subsection{Separate symmetry statuses}

The formalism distinguishes:
\begin{enumerate}
\item \textbf{local gauge symmetry}, implemented through
\cref{sec:gauge};
\item \textbf{exact global charges} of the regulated model, such as baryon
number and any exactly retained flavor charges;
\item \textbf{approximate internal symmetries}, such as isospin before
\(m_d-m_u\), QED, and nuclear breaking terms;
\item \textbf{front-form kinematics}, including exact kinematical generators
at the declared truncation;
\item \textbf{dynamical spacetime symmetries}, especially transverse
rotations, whose restoration is tested along the truncation tower;
\item \textbf{discrete maps} \(P,T,C\), whose light-front realization and
phase conventions are declared explicitly.
\end{enumerate}
These structures are related but do not form one exact direct-product global
group acting identically on every regulated object.

\subsection{Helicity and orbital angular momentum}

For a Fock amplitude in an intrinsic angular basis,
\begin{equation}
 \Lambda
 =\sum_{i=1}^{N_\nu}\lambda_i+L_z,
\qquad
 L_z=\sum_{j=1}^{N_\nu-1}\ell_{z,j},
\label{eq:jz-ledger}
\end{equation}
where only \(N_\nu-1\) intrinsic transverse momenta are independent.  In a
continuous momentum representation, the orbital generator is
\begin{equation}
 L_z
 =-i\sum_{j=1}^{N_\nu-1}
 \left(k_{jx}\frac{\partial}{\partial k_{jy}}
      -k_{jy}\frac{\partial}{\partial k_{jx}}\right)
\label{eq:oam-generator}
\end{equation}
with domain and boundary conditions inherited from the basis.  In a harmonic
or angular basis, \(L_z\) is a stored basis quantum number and
\cref{eq:oam-generator} defines its continuum meaning.

\begin{requirement}[OAM origin traceability]
\label{req:oam-traceability}
Every later TMD or GTMD harmonic attributed to orbital interference must
identify the pair of amplitude blocks whose \(L_z\), parton-helicity, and
target-helicity difference supplies the required transverse weight.  If
either block is disabled, the harmonic must vanish within numerical
tolerance.
\end{requirement}

\subsection{Equivariant couplings}

For a represented exact or controlled symmetry \(G\), an allowed coupling is
an intertwiner
\begin{equation}
 C:\bigotimes_{i=1}^nV_{R_i}\longrightarrow V_{R_{\rm out}},
 \qquad
 C D_{\rm in}(g)=D_{\rm out}(g)C.
\label{eq:general-intertwiner}
\end{equation}
The multiplicity space
\begin{equation}
 \Hom_G\!\left(\bigotimes_iV_{R_i},V_{R_{\rm out}}\right)
\label{eq:intertwiner-space}
\end{equation}
is part of the parameterization: independent invariant couplers may carry
independent interaction coefficients only when the Hamiltonian admits them.
Selection rules are implemented by absent blocks, not by fitting a large
penalty to forbidden matrix elements.

\begin{proposition}[Exact-charge selection rule]
\label{prop:selection-rule}
Let \(Q\) be self-adjoint, \(\comm{Q}{O}=0\), and let
\(Q|\nu,a\rangle=q_{\nu,a}|\nu,a\rangle\).  Then
\begin{equation}
 \langle\mu,b|O|\nu,a\rangle=0
 \quad\text{whenever}\quad q_{\mu,b}\ne q_{\nu,a}.
\label{eq:selection-rule}
\end{equation}
\end{proposition}

\begin{proof}
Taking the matrix element of \([Q,O]=0\) gives
\((q_{\mu,b}-q_{\nu,a})\langle\mu,b|O|\nu,a\rangle=0\).
\end{proof}

\subsection{Controlled isospin breaking}

A useful decomposition of the mass operator is
\begin{equation}
 \cM_r^2
 =\cM_{r,0}^2
 +\delta\cM_{r,m_d-m_u}^2
 +\delta\cM_{r,\rm QED}^2
 +\delta\cM_{r,\rm other}^2.
\label{eq:isospin-breaking-mass}
\end{equation}
The exact isospin limit is recovered by disabling the breaking operators, not
by assigning equal \(u\)- and \(d\)-flavor amplitudes inside a physical
proton.  Proton and neutron states remain distinct members of an isospin
multiplet, allowing exact charge-symmetry tests and controlled breaking.

\subsection{Rotational and Poincare defects}

Let \(G_\alpha^{(r)}\) be regulated generators evaluated on a selected
low-energy subspace \(\cS_r\).  Define
\begin{equation}
 \Delta_{\alpha\beta}^{(r)}
 =\left.
 \left(
 [G_\alpha^{(r)},G_\beta^{(r)}]
 -i f_{\alpha\beta}{}^\gamma G_\gamma^{(r)}
 \right)\right|_{\cS_r},
\label{eq:poincare-defect-operator}
\end{equation}
and a dimensionless defect
\begin{equation}
 \epsilon_{\alpha\beta}^{(r)}
 =\frac{\norm{\Delta_{\alpha\beta}^{(r)}}}
 {1+\sum_\gamma\abs{f_{\alpha\beta}{}^\gamma}
 \norm{G_\gamma^{(r)}}}.
\label{eq:poincare-defect}
\end{equation}
The operator norm may be replaced by a Frobenius or state-weighted norm only
if the choice is recorded.  Convergence requires the relevant defects to
decrease as the basis, Fock content, and counterterm basis are enlarged.
Degenerate \(J\)-multiplet masses and current angular conditions provide
independent rotational diagnostics.

\section{The regulated light-front mass operator}
\label{sec:mass-operator}

\subsection{Free invariant mass}

In sector \(\nu\) and the \(\PT=0\) frame, the free mass-squared kernel is
\begin{equation}
 \cM_{0,\nu}^{2(r)}
 (\{x_i,\bm k_{iT}\})
 =\sum_{i=1}^{N_\nu}
 \frac{m_{i,r}^2+\bm k_{iT}^2}{x_i},
\label{eq:free-invariant-mass}
\end{equation}
with masses in the declared Hamiltonian scheme.  Endpoint and ultraviolet
singularities are regulated by \(r\); a basis cutoff is not reinterpreted as
a physical constituent form factor unless it is matched as one.

\subsection{Operator decomposition}

At regulator level \(r\), write
\begin{equation}
 \cM_r^2
 =\cM_{0,r}^2
 +\cV_{\rm can,r}
 +\cV_{\rm eff,r}
 +\cV_{\rm ind,r}
 +\delta\cM_{\rm ct,r}^2.
\label{eq:mass-decomposition}
\end{equation}
The terms have different statuses:
\begin{description}[leftmargin=3.2cm,style=nextline]
\item[\(\cV_{\rm can,r}\)] Regulated canonical interactions obtained from
the selected light-front QCD Hamiltonian and constrained-field elimination.
\item[\(\cV_{\rm eff,r}\)] Explicit effective interactions, such as a
low-resolution confining or chiral operator, with power counting and
calibration conditions.
\item[\(\cV_{\rm ind,r}\)] Interactions induced by eliminating degrees of
freedom or by a similarity/unitary transformation.
\item[\(\delta\cM_{\rm ct,r}^2\)] Counterterms required by ultraviolet,
infrared, sector, symmetry-restoration, or truncation renormalization.
\end{description}
These classes must not be merged into an anonymous trainable matrix.

\subsection{Block structure}

Using the projectors of \cref{ax:orthogonal-sectors}, define
\begin{equation}
 \cM_{\nu\mu}^{2(r)}
 =P_\nu^{(r)}\cM_r^2P_\mu^{(r)},
\label{eq:mass-block}
\end{equation}
so that
\begin{equation}
 \cM_r^2
 =\sum_{\nu,\mu\in\cF_r}\cM_{\nu\mu}^{2(r)}.
\label{eq:mass-block-sum}
\end{equation}
For the minimal nucleon sectors,
\begin{equation}
 \cM_r^2\sim
 \begin{pmatrix}
 \cM_{qqq,qqq}^2 & V_{qqq,qqqg} & V_{qqq,qqqq\bar q} & \cdots\\
 V_{qqqg,qqq} & \cM_{qqqg,qqqg}^2 & V_{qqqg,qqqq\bar q} & \cdots\\
 V_{qqqq\bar q,qqq} & V_{qqqq\bar q,qqqg} &
 \cM_{qqqq\bar q,qqqq\bar q}^2 & \cdots\\
 \vdots&\vdots&\vdots&\ddots
 \end{pmatrix}.
\label{eq:nucleon-block-matrix}
\end{equation}
Each nonzero block must identify the generating operator, color and spin
intertwiner, momentum kernel, regulator, coefficient, adjoint partner, and
any subtraction or counterterm relation.

\subsection{Self-adjointness}

Let \(\cD_r\subset\Phys^{(r)}\) be a dense invariant domain in an infinite
regulated space, or the complete space in a finite matrix realization.

\begin{axiom}[Self-adjoint mass operator]
\label{ax:self-adjoint-mass}
The physical mass operator satisfies
\begin{equation}
 \cM_r^2=(\cM_r^2)^\dagger,
 \qquad
 \cD(\cM_r^2)=\cD((\cM_r^2)^\dagger).
\label{eq:self-adjoint-mass}
\end{equation}
\end{axiom}

\begin{proposition}[Block adjoint criterion]
\label{prop:block-adjoint}
On an orthogonal direct sum with a common invariant domain,
\cref{eq:self-adjoint-mass} implies
\begin{equation}
 (\cM_{\nu\mu}^{2(r)})^\dagger
 =\cM_{\mu\nu}^{2(r)}.
\label{eq:block-adjoint}
\end{equation}
Conversely, in a finite-dimensional regulated realization, self-adjoint
diagonal blocks together with \cref{eq:block-adjoint} imply that the full
block matrix is self-adjoint.
\end{proposition}

\begin{proof}
Using \(P_\nu^\dagger=P_\nu\),
\((P_\nu\cM^2P_\mu)^\dagger=P_\mu\cM^2P_\nu\).  The converse follows by
summing the block identities in a finite orthogonal direct sum.
\end{proof}

\subsection{Canonical interaction classes}

The formalism does not commit here to one regulator-specific kernel, but the
canonical QCD content must account for:
\begin{enumerate}
\item quark kinetic and mass terms;
\item transverse-gluon kinetic terms;
\item quark-gluon emission and absorption;
\item three- and four-gluon interactions;
\item instantaneous-fermion and instantaneous-gauge interactions generated
by constrained-field elimination;
\item regulator fields or subtraction terms when present.
\end{enumerate}
A canonical vertex has the typed form
\begin{equation}
 V_A^{(r)}:
 \Hil_{\mu,\rm phys}^{(r)}
 \longrightarrow
 \Hil_{\nu,\rm phys}^{(r)},
\qquad
 V_A^{(r)}=\cP_{\nu}^{\rm phys}K_A^{(r)}C_A S_A
 \cP_{\mu}^{\rm phys},
\label{eq:typed-vertex}
\end{equation}
where \(K_A\) is the momentum kernel, \(C_A\) the color intertwiner, and
\(S_A\) the helicity/spin tensor.  Their multiplication order and index
contractions are part of the operator definition.

\subsection{Effective and induced interactions}

An effective interaction is admissible only with a declaration
\begin{equation}
 \mathfrak v_A=
 (O_A,c_A,\Lambda_A,\nu_A,\mathsf{PC}_A,
 \bm\cC_A,\mathsf{replace}_A),
\label{eq:effective-interaction-record}
\end{equation}
containing the operator, coefficient, resolution scale, sectors, power
counting, calibration conditions, and replacement/subtraction semantics.  A
confining or chiral term may be central in a low-resolution model, but it
cannot be described as canonical QCD merely because it has the correct
quantum numbers.

\begin{principle}[Parameter locality]
\label{prin:parameter-locality}
A parameter belongs to a Hamiltonian term, current, matching coefficient,
state boundary condition, or explicit discrepancy operator.  It propagates
to every observable containing that object.  There is no parameter whose
primitive meaning is ``the normalization of named TMD \(F_A\).''
\end{principle}

\section{Eliminated sectors, effective Hamiltonians, and induced operators}
\label{sec:feshbach}

\subsection{Exact projected equations}

Let \(P\) project onto a retained subspace and \(Q=1-P\).  For
\begin{equation}
 H|\Psi\rangle=E|\Psi\rangle,
\label{eq:generic-eigenproblem}
\end{equation}
the coupled equations are
\begin{align}
 (PHP-E)P|\Psi\rangle+PHQ\,Q|\Psi\rangle&=0,
\label{eq:p-equation}\\
 QHP\,P|\Psi\rangle+(QHQ-E)Q|\Psi\rangle&=0.
\label{eq:q-equation}
\end{align}
When the resolvent exists,
\begin{equation}
 Q|\Psi\rangle
 =\omega(E)P|\Psi\rangle,
 \qquad
 \omega(E)=(E-QHQ)^{-1}QHP.
\label{eq:wave-operator}
\end{equation}
The exact energy-dependent retained-space Hamiltonian is
\begin{equation}
 H_{\rm eff}(E)
 =PHP+PHQ(E-QHQ)^{-1}QHP.
\label{eq:feshbach-hamiltonian}
\end{equation}
For a bound state below eliminated-sector thresholds and a self-adjoint
starting Hamiltonian, the resolvent can yield a Hermitian effective operator.
Above thresholds the boundary prescription and possible non-Hermiticity must
be declared.

\subsection{Eliminating a sector also changes operators}

For any operator \(O\), matrix elements between exact states can be written
in the retained space as
\begin{equation}
 \langle\Psi'|O|\Psi\rangle
 =\langle\psi'_P|
 O_{\rm eff}(E',E)|\psi_P\rangle,
\label{eq:effective-operator-matrix}
\end{equation}
where \(|\psi_P\rangle=P|\Psi\rangle\) and
\begin{equation}
 O_{\rm eff}(E',E)
 =P\bigl(1+\omega^\dagger(E')\bigr)
 O\bigl(1+\omega(E)\bigr)P.
\label{eq:effective-operator}
\end{equation}
The retained vectors carry the norm kernel
\begin{equation}
 N(E)=P\bigl(1+\omega^\dagger(E)\omega(E)\bigr)P.
\label{eq:norm-kernel}
\end{equation}
An energy-independent Hermitian effective theory may instead be obtained by a
specified similarity or unitary transformation, but the transformed currents
and composite operators must accompany the transformed Hamiltonian.

\begin{theorem}[No Hamiltonian-only elimination]
\label{thm:no-hamiltonian-only-elimination}
If an eliminated sector contributes to \(\omega(E)\) and an observable
operator \(O\) connects to that sector, replacing \(H\) by
\(H_{\rm eff}\) while retaining \(POP\) alone does not, in general,
preserve the exact matrix element.
\end{theorem}

\begin{proof}
The exact retained-space operator is
\cref{eq:effective-operator}.  Unless
\(PO\omega=\omega^\dagger OP=\omega^\dagger O\omega=0\), it differs from
\(POP\).  Therefore the Hamiltonian reduction and operator reduction must be
matched together.
\end{proof}

This theorem is central to the later nuclear hierarchy: integrating out an
\(NN\pi\) sector induces both an effective \(NN\) interaction and pion or
exchange-current contributions to partonic and electromagnetic operators.

\subsection{Double-counting implication}

If physics from an eliminated sector is represented by
\(\cV_{\rm ind}\) or \(O_{\rm eff}\), an explicit version of the same
sector may not be added again without a subtraction map.  The provenance
record must contain a two-cell of the form
\begin{equation}
 \text{explicit sector path}
 \quad\Longleftrightarrow\quad
 \text{induced-operator path}+\text{subtraction/remainder}.
\label{eq:replacement-two-cell}
\end{equation}
A closed transition cycle alone is not proof of double counting; the
replacement identity in \cref{eq:replacement-two-cell} is the operative
evidence.

\section{Renormalization datum and directed truncation system}
\label{sec:renormalization}

\subsection{Renormalization conditions}

Let the regulated coefficient vector be \(\bm c_r\).  Renormalization
conditions are functionals
\begin{equation}
 \cC_i[\mathfrak T_r]=c_i^{\rm target},
 \qquad i=1,\ldots,N_{\rm ren},
\label{eq:renormalization-conditions}
\end{equation}
which may fix masses, charges, selected form factors, symmetry-restoration
conditions, or EFT low-energy constants.  The full record is
\begin{equation}
 \mathfrak R_r=
 (\cR_r,\bm c_r,\bm\cC_r,
 \mathfrak s_r,\mu_r,\mathsf{solve}_r,
 \operatorname{Cov}_{\rm ren,r}).
\label{eq:renormalization-record}
\end{equation}
It includes the solve procedure and covariance or numerical sensitivity of
the fitted conditions.

\subsection{Sector-dependent renormalization}

Fock truncation can require bare masses, couplings, or counterterms that
depend on the number of sectors available above a vertex.  Denote such a
coefficient by
\begin{equation}
 c_{A}^{(r;\nu\leftarrow\mu)}.
\label{eq:sector-dependent-coefficient}
\end{equation}
Sector dependence is admissible only when:
\begin{enumerate}
\item the dependence follows a declared truncation-renormalization scheme;
\item common physical renormalization conditions are recovered as the sector
space is enlarged;
\item Ward or Slavnov--Taylor constraints are tested where applicable;
\item the dependence is not used to fit unrelated observables independently;
\item the limiting or matching relation among sector coefficients is stored.
\end{enumerate}

\begin{requirement}[Counterterm completeness study]
\label{req:counterterm-completeness}
For each enlargement of regulator or Fock content, the operator basis of
counterterms must be checked for closure under the retained symmetries and
power counting.  Stability obtained by refitting an incomplete arbitrary
matrix is not a renormalization result.
\end{requirement}

\subsection{Directed family of calculations}

Let \((\mathfrak R,\preceq)\) be a directed set of regulator/truncation
labels.  The relation \(r\preceq r'\) means that \(r'\) is at least as
resolved along specified axes; it does not require
\(\Hil_r\subseteq\Hil_{r'}\).  The system may contain:
\begin{equation}
 I_{r\to r'}:\Hil_r\longrightarrow\Hil_{r'},
 \qquad
 R_{r'\to r}:\Alg_{r'}\longrightarrow\Alg_r,
\label{eq:comparison-maps}
\end{equation}
where \(I\) is a prolongation/comparison map and \(R\) is an effective
operator or coarse-graining map.  Their existence and accuracy are separate
claims.

A common-reference matching map is often cleaner:
\begin{equation}
 \cZ_{r\to\star}:\Alg_r\longrightarrow\Alg_\star,
\label{eq:common-reference-matching}
\end{equation}
where \(\star\) is a declared renormalized low-energy or QCD operator scheme.
Matched observables are
\begin{equation}
 \cO_i^{[\star]}(r)
 =\langle\Psi_r|
 \cZ_{r\to\star}[\cO_{i,r}]
 |\Psi_r\rangle.
\label{eq:matched-observable}
\end{equation}

\begin{definition}[Observable convergence]
\label{def:observable-convergence}
A family is convergent for an observable set \(\mathfrak O\) if there exist
limits \(\cO_{i,\infty}^{[\star]}\) such that
\begin{equation}
 \lim_{r\to\infty}
 \cO_i^{[\star]}(r)
 =\cO_{i,\infty}^{[\star]}
 \qquad\text{for all }i\in\mathfrak O,
\label{eq:observable-convergence}
\end{equation}
with the directed limit taken along every declared required truncation axis.
\end{definition}

No analogous requirement is imposed on bare wave-function coefficients from
different bases or schemes unless a specific isometric comparison map is
provided.

\subsection{Truncation axes and error model}

At minimum, keep distinct axes for:
\begin{equation}
 \nu_{\rm trunc}=
 (K,N_{\max},\cF,\Lambda_{\rm UV},\Lambda_{\rm IR},
 r_{\rm network},n_{\rm quad},\epsilon_{\rm eig},\mathsf{ct}).
\label{eq:truncation-axes}
\end{equation}
Here \(r_{\rm network}\) is a tensor-network or factorization rank when used,
and \(\mathsf{ct}\) labels the counterterm basis.  A controlled asymptotic
model may take the form
\begin{equation}
 \cO_i(r)
 =\cO_{i,\infty}
 +\sum_a A_{ia}\,\varepsilon_a(r)^{p_{ia}}
 +\sum_{a<b}A_{iab}\varepsilon_a(r)^{p_{ia}}
 \varepsilon_b(r)^{p_{ib}}+\cdots,
\label{eq:convergence-model}
\end{equation}
with priors or ranges for unknown exponents documented.  The convergence fit
must be tested on at least one withheld resolution level whenever enough
levels exist.

\begin{principle}[No uncertainty substitution]
\label{prin:no-uncertainty-substitution}
Posterior parameter uncertainty, Fock truncation, basis truncation,
regulator dependence, matching uncertainty, symmetry defects, and numerical
error are separate random or bounded quantities.  Convergence in one cannot
be used to hide the absence of another.
\end{principle}

\subsection{Projective consistency of operators}

For operators matched from a finer to a coarser level, require
\begin{equation}
 R_{r'\to r}[\cO_{A,r'}]
 =\cO_{A,r}+\delta\cO_{A;r'\to r},
\label{eq:operator-projective-consistency}
\end{equation}
where the remainder is bounded on a declared state class:
\begin{equation}
 \sup_{\psi\in\mathfrak S_r}
 \frac{\abs{\langle\psi|\delta\cO_{A;r'\to r}|\psi\rangle}}
 {\cN_A(\psi)}
 \le \epsilon_{A}(r',r).
\label{eq:operator-remainder-bound}
\end{equation}
The bound, not the visual similarity of two curves, is the formal convergence
statement.

\section{Eigenstates, phase conventions, and conserved ledgers}
\label{sec:states}

\subsection{Regulated eigenproblem}

At each regulator level, solve
\begin{equation}
 \cM_r^2|\Psi_{h,n};P^+,\PT,\Lambda\rangle_r
 =M_{h,n}^{2(r)}
 |\Psi_{h,n};P^+,\PT,\Lambda\rangle_r,
\label{eq:regulated-eigenproblem}
\end{equation}
with
\begin{equation}
 \langle\Psi_{h,n};P',\Lambda'|
 \Psi_{h,m};P,\Lambda\rangle_r
 =2P^+(2\pi)^3\delta^{(3)}_{\LF}(P'-P)
 \delta_{nm}\delta_{\Lambda'\Lambda}.
\label{eq:hadron-state-normalization}
\end{equation}
After the overall momentum delta is removed, the internal states have unit
norm as in \cref{eq:fock-normalization}.  The eigensolver must return a
residual
\begin{equation}
 \epsilon_{\rm eig}
 =\frac{\norm{\cM_r^2|\Psi\rangle-M_r^2|\Psi\rangle}}
 {\max(1,\norm{\cM_r^2|\Psi\rangle},|M_r^2|)}
\label{eq:eigensolver-residual}
\end{equation}
and an orthogonality defect for every retained state cluster.

\subsection{Phase and degenerate-subspace conventions}

An eigenvector is defined only up to phase, and a degenerate eigenspace up to
unitary rotation.  Since GTMD overlaps are phase sensitive, each state family
requires a transport convention.  For an isolated state along a parameter or
truncation path \(r\to r'\), choose the phase so that
\begin{equation}
 \langle\Psi_h^{(r')}|I_{r\to r'}|\Psi_h^{(r)}\rangle
 \in\mathbb R_{\ge0}
\label{eq:phase-transport}
\end{equation}
when the overlap is nonzero.  For a nearly degenerate multiplet, align the
subspaces by the unitary factor in the polar decomposition of the overlap
matrix rather than matching individual vectors greedily.

\begin{requirement}[State identity]
\label{req:state-identity}
A serialized state must contain: hadron and excitation identity, total
momentum convention, \(J^z\)/helicity, regulator and basis identifiers,
sector ordering, inner-product metric, phase convention, eigenvalue,
residual, solver configuration, source commit, and content hash.
\end{requirement}

\subsection{Exact-charge ledgers}

Let \(Q_A\) be an exact conserved additive charge with constituent eigenvalue
\(q_{A,i}\).  In each sector,
\begin{equation}
 Q_A|\nu,\bm\alpha\rangle
 =\left(\sum_{i\in\nu}q_{A,i}\right)
 |\nu,\bm\alpha\rangle.
\label{eq:additive-charge}
\end{equation}
For a physical state of charge \(q_A^{(h)}\), require
\begin{equation}
 \epsilon_{Q_A}
 =\norm{(Q_A-q_A^{(h)})|\Psi_h\rangle}
\label{eq:charge-defect}
\end{equation}
to vanish exactly in a block-diagonal finite representation or to remain
below its declared regulator tolerance.  Ledgers are required for baryon
number, electric charge, exact flavor charges, and total color.

\subsection{Plus-momentum ledger}

In every basis configuration on \(\cM_N\),
\begin{equation}
 \sum_i x_i=1.
\label{eq:configuration-momentum-ledger}
\end{equation}
The integrated constituent momentum fractions satisfy
\begin{equation}
 \sum_{a\in\mathsf A}
 \langle x\rangle_{a,r}
 =\sum_{\nu,\bm\alpha}
 \int[\dd\Gamma_{N_\nu}]
 |\psi_{\nu,\bm\alpha}^{(r)}|^2
 \sum_{i\in\nu}x_i
 =1.
\label{eq:state-momentum-ledger}
\end{equation}
This identity is regulator-exact for a complete normalized basis state.  A
later partonic PDF momentum sum rule additionally requires matched composite
operators and is not inferred from \cref{eq:state-momentum-ledger} alone.

\subsection{Helicity and orbital ledger}

For a \(J^z=\Lambda\) eigenstate,
\begin{equation}
 \Lambda
 =\sum_{\nu,\bm\alpha}
 \int[\dd\Gamma_{N_\nu}]
 |\psi_{\nu,\bm\alpha}|^2
 \left(\sum_i\lambda_i+L_z\right)_{\nu,\bm\alpha}
\label{eq:state-jz-ledger}
\end{equation}
when the basis diagonalizes the indicated operators.  In a non-diagonal
basis, the right side is the corresponding expectation value.  This is a
state-space identity for canonical front-form \(J^z\).  The later
renormalized quark/gluon spin decomposition requires scheme-defined local or
nonlocal operators and is a separate statement.

\subsection{Eigensolver admissibility}

Dense diagonalization, Lanczos, Davidson, tensor-network variational methods,
quantum eigensolvers, or other algorithms are admissible only when they solve
the same typed operator.  For every reported state, record:
\begin{itemize}
\item matrix/operator hash and domain;
\item symmetry block and preconditioner;
\item starting vectors and random seeds;
\item stopping criteria and residual history;
\item orthogonality and degeneracy handling;
\item comparison with exact diagonalization on every tractable smaller
space;
\item stability under increased numerical precision and solver parameters.
\end{itemize}
A variational energy alone is insufficient if the wavefunction will be used
in interference observables.

\section{Currents and local operators consistent with the dynamics}
\label{sec:currents}

\subsection{Hamiltonian-current pair}

A microscopic state cannot be combined with an unrelated current while
retaining a first-principles claim.  At regulator level \(r\), write
\begin{equation}
 J_r^\mu
 =J_{1\rm b,r}^\mu
 +J_{\rm exch,r}^\mu
 +J_{\rm ind,r}^\mu
 +\delta J_{\rm ct,r}^\mu,
\label{eq:current-decomposition}
\end{equation}
where the exchange and induced terms follow from the retained interactions,
sector elimination, and matching.  Analogous decompositions apply to the
axial current, energy-momentum tensor, and the bilocal operators of later
volumes.

\begin{definition}[Consistent operator family]
\label{def:consistent-operator-family}
A family \(\{\cM_r^2,\cO_{A,r}\}\) is consistent when every operator is
constructed or matched in the same regulator, basis, gauge, zero-mode, and
renormalization datum, including induced terms generated by any
transformation used to define \(\cM_r^2\).
\end{definition}

\subsection{Ward--Takahashi and charge conditions}

For a conserved electromagnetic current,
\begin{equation}
 q_\mu\langle P',\Lambda'|J_r^\mu(0)|P,\Lambda\rangle
 =\delta_{\rm Ward}^{(r)}(P',P;\Lambda',\Lambda),
\label{eq:ward-identity}
\end{equation}
where the defect must converge to zero and is not removed by rescaling the
final form factor.  Charge normalization requires
\begin{equation}
 \frac{1}{2P^+}
 \langle P,\Lambda|J_r^+(0)|P,\Lambda\rangle
 =Q_h+\delta_Q^{(r)}.
\label{eq:charge-normalization}
\end{equation}
A charge counterterm may be fixed by \cref{eq:charge-normalization}; the
remaining \(q^2\) and helicity dependence is then predictive within the
truncation and current basis.

\subsection{Spin-1 angular condition}

When the eventual deuteron state is evaluated in a Drell--Yan frame
\(q^+=0\), four independent plus-current helicity matrix elements are often
used to extract three physical elastic form factors.  In a common helicity
phase convention their redundancy gives the angular condition
\begin{equation}
 \Delta_{\rm ang}(Q^2)
 =(1+2\eta)I^+_{11}
 +I^+_{1,-1}
 -\sqrt{8\eta}\,I^+_{10}
 -I^+_{00}=0,
 \qquad
 \eta=\frac{Q^2}{4M_D^2}.
\label{eq:spin1-angular-condition}
\end{equation}
Other phase conventions transform the displayed signs together with the
helicity states.  The implementation therefore stores the helicity phase
adapter and evaluates an invariantly derived condition, not a hard-coded
sign string.

A violation of \cref{eq:spin1-angular-condition} diagnoses missing dynamical
rotations, omitted current components, basis truncation, or inconsistent
matching.  Choosing one three-element subset of helicity amplitudes does not
make the violation disappear.

\subsection{Local operator renormalization}

Composite operators require their own renormalization and can mix:
\begin{equation}
 \cO_{A}^{[\star]}(\mu)
 =\sum_B Z_{AB}^{r\to\star}(\mu)
 \cO_{B,r}+\delta\cO_{A,r}^{\rm power}.
\label{eq:local-operator-renormalization}
\end{equation}
The matching matrix must preserve exact charges and declared Ward identities.
For energy-momentum operators, quark and gluon components can mix even though
the conserved total tensor is protected.  Bare constituent momentum ledgers
must therefore not be confused with matched partonic operator fractions.

\subsection{Minimum local-current validation set}

Before a state enters Volume~II, evaluate where applicable:
\begin{enumerate}
\item electric charge and flavor charges;
\item magnetic moment or magnetic form-factor normalization;
\item axial charge;
\item low moments of the energy-momentum tensor;
\item current conservation and Ward identities;
\item rotational/angular conditions for spin \(\ge1\);
\item stability under basis, Fock, regulator, current-basis, and numerical
enlargement.
\end{enumerate}
The purpose is not to fit every item.  A restricted subset fixes the
renormalization datum; the remainder tests the shared state-current pair.

\section{Typed algebra and the C1 convention/operator spine}
\label{sec:c1-spine}

\subsection{Purpose of C1}

The Stage~0 audit found no accepted numerical convention conflict.  C1 is
therefore an adapter migration, not a physics rewrite.  It introduces a
fail-closed identity spine around the existing model so that the accepted
arrays can no longer be composed through generic coordinate, rank, path, or
map aliases.

\begin{principle}[Adapter-first migration]
\label{prin:adapter-first}
A legacy object is wrapped by an immutable typed view.  The adapter may
expose existing values and reproduce existing formulas; it may not change
values, reorder rows, insert defaults that affect physics, or rewrite an
authoritative artifact.  Native typed objects replace adapters only after
byte-identical regression is demonstrated.
\end{principle}

\subsection{Coordinate specification}

Define a coordinate record
\begin{equation}
\begin{aligned}
 \mathsf{CoordinateSpec}=(&\texttt{kind},d,\texttt{units},\texttt{frame},\\
 &\texttt{role},\texttt{conjugate},\texttt{sign},\texttt{measure}).
\end{aligned}
\label{eq:coordinate-spec}
\end{equation}
The enumeration \texttt{kind} contains at least the eight kinds of
\cref{sec:typed-transverse}.  A vector value is a pair
\begin{equation}
 \mathsf{TypedVector}=(\bm v,\mathsf{CoordinateSpec}).
\label{eq:typed-vector}
\end{equation}
Dot products and Fourier kernels require compatible dimensions, units,
frames, and a declared conjugacy relation.

\begin{proposition}[No implicit Wigner--TMD transform]
\label{prop:no-implicit-wigner-tmd}
Under \cref{ax:no-transverse-alias}, a transform
\(\DT\leftrightarrow\bD\) cannot be substituted for
\(\kT\leftrightarrow\bM\) by structural typing, even though both use a
two-dimensional Fourier kernel.
\end{proposition}

\begin{proof}
The transform domain and codomain carry different \texttt{kind} labels and
different physical roles.  The composition predicate in
\cref{eq:composition-predicate} rejects unequal types unless an explicit
adapter is registered.
\end{proof}

\subsection{Transverse-rank and mass-normalization specification}

Define
\begin{equation}
\begin{aligned}
 \mathsf{RankSpec}=(&\texttt{weight},\texttt{tensor\_basis},
 \texttt{k\_power},\texttt{b\_power},\\
 &M_{\rm ref},\texttt{bessel\_order},
 \texttt{fourier\_phase},\\
 &\texttt{coefficient\_role}).
\end{aligned}
\label{eq:rank-software-spec}
\end{equation}
Here \texttt{coefficient\_role} distinguishes a scalar coefficient from the
fully contracted physical modulation.  The rank may be zero even when the
named observable contains a target tensor; it is the transverse harmonic
weight of the coefficient/tensor pair that controls the Bessel kernel.

A transform is admissible only if
\begin{equation}
 \texttt{bessel\_order}=|\texttt{weight}|
\label{eq:bessel-rank-condition}
\end{equation}
and the stored powers make the forward and inverse dimensions reciprocal.
C1 need not replace the accepted transform formulas; it must make their
metadata explicit and test that a wrong rank, mass, or phase is rejected.

\subsection{Sector identity}

Define
\begin{equation}
\begin{aligned}
 \mathsf{SectorId}=(&\texttt{resolution},\bm n_q,\bm n_{\bar q},n_g,
 \bm Q_{\rm exact},J^z,\\
 &\texttt{parity\_class},\texttt{color\_status},
 \texttt{basis\_id}).
\end{aligned}
\label{eq:sector-id}
\end{equation}
The object is immutable and hashable.  Its exact charges are checked against
occupations at construction.  A state tensor may not silently concatenate
rows from different \texttt{SectorId} values merely because their numerical
shapes agree.

\subsection{Wilson path identity}

A bare path identity is
\begin{equation}
\begin{aligned}
 \mathsf{WilsonPathId}=(&\texttt{endpoints},\texttt{segments},
 \texttt{orientation},\texttt{closure},\\
 &\texttt{rapidity\_direction},\texttt{representation},
 \texttt{boundary\_class}).
\end{aligned}
\label{eq:wilson-path-id}
\end{equation}
For a quark bilocal there is one represented path.  For a gluon correlator,
store an ordered pair
\begin{equation}
 \mathsf{GluonLinkId}
 =(\mathsf{WilsonPathId}_1,\mathsf{WilsonPathId}_2,
 \texttt{color\_class}),
\label{eq:gluon-link-id}
\end{equation}
where \texttt{color\_class} is one of
\(\{\texttt{F\_TYPE},\texttt{D\_TYPE},
\texttt{NOT\_APPLICABLE},\texttt{UNSPECIFIED}\}\).  Production evaluation
refuses \texttt{UNSPECIFIED} whenever the observable depends on the class.

\subsection{Decorated operator identity}

A complete operator identity is
\begin{equation}
\begin{aligned}
 \mathsf{DecoratedOperatorId}=\bigl(&
 \texttt{name},\texttt{species},\texttt{flavor},\texttt{projection},\\
 &\texttt{domain},\texttt{codomain},\texttt{momentum\_fibers},\\
 &\texttt{coordinate\_specs},\mathsf{RankSpec},
 \texttt{wilson\_identity},\\
 &\texttt{color\_representation},\texttt{uv\_regulator},\\
 &\texttt{rapidity\_regulator},\texttt{soft\_prescription},\mu,\zeta,\\
 &\texttt{scheme},\texttt{normalization},\texttt{version}\bigr).
\end{aligned}
\label{eq:decorated-operator-id}
\end{equation}
Fields not relevant to a local or bare operator carry an explicit
\texttt{NOT\_APPLICABLE}; they are not silently omitted.  An unknown field is
\texttt{UNSPECIFIED}, which is distinct from not applicable and fails any
operation that requires it.

\subsection{Map classes}

Every map carries one of the disjoint classes
\begin{equation}
 \mathsf{MapClass}
 \in\{\texttt{AMP},\texttt{DENS},\texttt{MATCH},
 \texttt{RED},\texttt{PROC}\}.
\label{eq:map-class-enum}
\end{equation}
Their meanings are:
\begin{longtable}{L{0.14\textwidth}L{0.28\textwidth}L{0.49\textwidth}}
\toprule
Class & Primitive objects & Permitted primitive meaning\\
\midrule
\endhead
\texttt{AMP} & Regulated Hilbert spaces and amplitude tensors & Hamiltonian blocks, coherent vertices, current insertions, state embeddings, adjoints.\\
\texttt{DENS} & Operator algebras and density objects & Partial traces, selected incoherent outcomes, positive or completely positive reductions after amplitudes are formed.\\
\texttt{MATCH} & Regulator-, resolution-, and scheme-indexed operator spaces & Renormalization, EFT matching, similarity transformations, soft subtraction, basis/truncation comparison.\\
\texttt{RED} & GTMD, TMD, GPD, current, and Wigner spaces & Typed limits, integrals, Fourier transforms, projections, and moments.\\
\texttt{PROC} & Factorized tensors and measured cross sections & Hard/soft/fragmentation/color assembly, measurement maps, factorization status, and \(W+Y\).\\
\bottomrule
\end{longtable}

\subsection{Fail-closed composition}

For typed maps
\begin{equation}
 f:(A,\tau_A)\to(B,\tau_B),
 \qquad
 g:(C,\tau_C)\to(D,\tau_D),
\label{eq:typed-maps}
\end{equation}
define the composition predicate
\begin{equation}
 \mathsf{CanCompose}(g,f)
 \Longleftrightarrow
 \bigl[(B,\tau_B)=(C,\tau_C)\bigr]
 \ \lor\ 
 \bigl[\exists\,\mathsf{Adapter}: (B,\tau_B)\to(C,\tau_C)\bigr].
\label{eq:composition-predicate}
\end{equation}
An adapter is itself a versioned typed map with a formula, provenance,
validation set, and loss/remainder declaration.  Shape compatibility is not
part of \(\mathsf{CanCompose}\).

\begin{theorem}[Fail-closed type safety]
\label{thm:fail-closed}
If all executable compositions are mediated by
\cref{eq:composition-predicate}, then no computation can compose mismatched
coordinate, scheme, operator, sector, or map-class identities without an
explicit registered adapter.
\end{theorem}

\begin{proof}
A direct composition is defined only for identical endpoint objects and
metadata.  Every nonidentical endpoint requires an adapter, which is itself
an explicit graph edge.  Therefore an unregistered mismatch has no defined
composition and fails before numerical evaluation.
\end{proof}

\subsection{C1 injected mismatch tests}

C1 must deliberately demonstrate rejection of at least:
\begin{enumerate}
\item \(\bD\) passed to a \(\bM\) transform;
\item a \(J_0\) transform applied to a rank-one or rank-two coefficient;
\item an absent reference mass where the coefficient convention requires one;
\item future and past link classes treated as identical;
\item a gluon \(f\)-type object substituted for a \(d\)-type object;
\item an \(\Amp\) map composed directly as a \(\Dens\) map;
\item a scheme-changing operation executed without a \(\Match\) adapter;
\item incompatible sector identities concatenated by array shape;
\item \texttt{UNSPECIFIED} operator metadata entering a production
observable that requires the field;
\item an adapter whose declared loss/remainder is missing.
\end{enumerate}
These are positive tests of refusal behavior, not merely unit tests of happy
paths.

\section{Computational realization and data contracts}
\label{sec:software}

\subsection{Core objects for Volume I}

\begin{longtable}{L{0.25\textwidth}L{0.65\textwidth}}
\toprule
Object & Formal responsibility\\
\midrule
\endhead
\texttt{CoordinateSpec}, \texttt{TypedVector} & Distinct transverse and longitudinal variable roles, units, frame, transform conjugacy, and measure.\\
\texttt{RankSpec} & Angular weight, tensor basis, mass normalization, extracted powers, Bessel order, and Fourier phase.\\
\texttt{RegulatorSpec} & Complete label \(r\) from \cref{eq:regulator-label}, including gauge and zero-mode policies.\\
\texttt{BasisSpec} & Basis family, mode ordering, normalization metric, quadrature/isometry, cutoffs, and serialization.\\
\texttt{SectorId} & Flavor-resolved occupations, exact charges, \(J^z\), color status, resolution, and basis identity.\\
\texttt{SectorSpace} & Kinematical and physical sector spaces, statistics, projectors, and support constraints.\\
\texttt{Intertwiner} & Sparse invariant tensor with group/gauge action, domain, codomain, multiplicity, and normalization tests.\\
\texttt{LFMassOperator} & Self-adjoint block operator with canonical/effective/induced/counterterm provenance and matrix-free or sparse application.\\
\texttt{FockState} & Normalized amplitudes, sector blocks, phase convention, eigenvalue, residual, and state identity.\\
\texttt{RenormalizationDatum} & Coefficients, conditions, scheme, solve record, covariance, and sector-dependent relations.\\
\texttt{CurrentOperator} & Local operator family matched consistently to the mass operator, including induced and counterterm pieces.\\
\texttt{TruncationLabel} and \texttt{TruncationTower} & Directed resolution levels, comparison/matching maps, convergence models, and required axes.\\
\texttt{ConvergenceManifest} & Matched observables, errors, withheld levels, fit diagnostics, symmetry defects, and acceptance outcome.\\
\texttt{WilsonPathId} and \texttt{DecoratedOperatorId} & C1 identity spine required now; dynamical Wilson evaluation belongs to Volume III.\\
\texttt{TypedMap} and \texttt{AdapterRegistry} & Map-class identity and fail-closed composition.\\
\bottomrule
\end{longtable}

\subsection{Sparse and tensor-network representation}

A tensor network may represent the state and operator contractions sparsely.
An edge label has at least
\begin{equation}
 e=(\nu,a,R_c,R_f,\lambda,L_z,J^z,
 \alpha,r,\mathfrak r,\mathfrak s),
\label{eq:network-edge}
\end{equation}
where every entry is physical or numerical metadata.  Nodes represent
Hamiltonian vertices, intertwiners, basis transforms, state tensors, currents,
matching maps, or controlled traces.  Block sparsity is generated from exact
selection rules and support, not from post hoc magnitude pruning.

A network bond or factorization rank is one truncation axis.  Results must be
reported against that rank together with basis, Fock, regulator, and
quadrature convergence.  A low-rank result is not accepted merely because it
fits the calibration data.

\subsection{Serialization and provenance}

Every durable object contains:
\begin{itemize}
\item a semantic type and schema version;
\item content hash and source commit;
\item dependency hashes;
\item units, normalization, basis, regulator, and scheme;
\item exact/model/numerical status labels;
\item creation command or callable provenance;
\item replacement, exclusion, and compatibility relations;
\item validation report reference.
\end{itemize}
Authoritative numerical data are content-addressed.  A filename is a human
alias, not identity.

\subsection{Matrix-free operator application}

Large calculations may avoid materializing \(\cM_r^2\).  A matrix-free
operator is acceptable if it exposes:
\begin{equation}
 \mathsf{apply}(v),\quad
 \mathsf{apply\_adjoint}(v),\quad
 \mathsf{domain},\quad
 \mathsf{codomain},\quad
 \mathsf{block\_map},\quad
 \mathsf{hash}.
\label{eq:matrix-free-interface}
\end{equation}
Hermiticity is tested by randomized bilinear forms,
\begin{equation}
 \epsilon_H(u,v)
 =\frac{\abs{\langle u|\cM^2v\rangle
 -\langle\cM^2u|v\rangle}}
 {1+\norm{u}\norm{\cM^2v}+\norm{\cM^2u}\norm{v}},
\label{eq:random-hermiticity-test}
\end{equation}
and by exact block comparisons wherever matrices are tractable.

\section{Formal propositions and required property tests}
\label{sec:tests}

\subsection{Structural propositions}

The following results are part of the formal contract already established:
\begin{enumerate}
\item kinematical boost invariance of intrinsic variables,
\cref{prop:intrinsic-boost};
\item nonnegative sector norm closure, \cref{prop:norm-ledger};
\item preservation of color singlets by invariant operators,
\cref{prop:singlet-preservation};
\item exact-charge block selection, \cref{prop:selection-rule};
\item block adjoint criterion, \cref{prop:block-adjoint};
\item induced-operator necessity under sector elimination,
\cref{thm:no-hamiltonian-only-elimination};
\item no implicit Wigner/TMD coordinate substitution,
\cref{prop:no-implicit-wigner-tmd};
\item fail-closed type safety, \cref{thm:fail-closed}.
\end{enumerate}
They should be represented directly as unit, property, or schema tests rather
than only cited in documentation.

\subsection{Volume I validation matrix}

\begin{longtable}{L{0.07\textwidth}L{0.35\textwidth}L{0.48\textwidth}}
\toprule
ID & Property & Minimum test\\
\midrule
\endhead
V1.1 & Light-front convention & Round-trip conversions reproduce \cref{eq:lf-coordinates,eq:lf-invariant}; deliberately altered factor-of-two convention is rejected.\\
V1.2 & Typed transverse coordinates & Every illegal pair in \cref{ax:no-transverse-alias} fails without an adapter.\\
V1.3 & Rank convention & Correct \(J_{|m|}\), mass powers, dimensions, and Fourier phase are stored; injected wrong metadata fail.\\
V1.4 & Basis metric & Orthonormality or declared Gram matrix is reproduced within tolerance; coefficient-only states are rejected.\\
V1.5 & Statistics & Exchange of identical fermions/bosons gives the required sign; factorial normalization closes once.\\
V1.6 & Sector grading & Projectors satisfy \cref{eq:sector-projector-algebra}; occupation-derived charges agree with \texttt{SectorId}.\\
V1.7 & Gauge/color & Singlet, Gauss/residual, and Ward defects are reported separately; invariant vertices preserve singlets.\\
V1.8 & Intertwiners & Each stored intertwiner commutes with represented exact generators within tolerance.\\
V1.9 & Mass operator & Diagonal blocks are self-adjoint and off-diagonal blocks satisfy \cref{eq:block-adjoint}; randomized bilinear tests pass.\\
V1.10 & Selection rules & Forbidden blocks are structurally absent; injected forbidden matrix elements are detected.\\
V1.11 & State solve & Eigen residual, orthogonality, phase/subspace tracking, and solver reproducibility pass.\\
V1.12 & Ledgers & Norm, exact charges, plus momentum, and \(J^z\) close with separate tolerances.\\
V1.13 & Sector elimination & A solvable benchmark reproduces the full-space energy and a nontrivial operator matrix element only when induced operators are included.\\
V1.14 & Renormalization & Conditions are satisfied and regulator variation is reported after refitting; coefficients and conditions are versioned.\\
V1.15 & Directed convergence & At least one matched observable is fitted across multiple levels and tested on a withheld level.\\
V1.16 & Current consistency & Charge, Ward identity, and where applicable angular condition are reported without after-the-fact rescaling.\\
V1.17 & Type composition & All map-class, scheme, path, sector, and coordinate mismatch injections fail closed.\\
V1.18 & Regression & C1 passes the full C0 suite and preserves all eight authoritative hashes byte for byte.\\
\bottomrule
\end{longtable}

\subsection{Analytic benchmark ladder}

Before a full QCD nucleon solve, validate the architecture on increasing
models:
\begin{enumerate}
\item a finite two-sector Hermitian matrix with an exact Feshbach reduction;
\item a scalar light-front two-body bound-state kernel with known
normalization and form factor;
\item a fermion-boson model testing helicity, sector-dependent
renormalization, and a Ward identity;
\item a color toy model with \(q\bar q\), \(qqq\), and \(qqqg\) singlet
intertwiners;
\item a small BLFQ-like basis block benchmarked by dense and sparse
solvers;
\item the selected microscopic nucleon pilot.
\end{enumerate}
The toy models test architecture; they do not count as evidence that the QCD
interaction has been solved.

\section{Contract for passage to Volume II}
\label{sec:volume2-contract}

\subsection{State bundle supplied to the GTMD layer}

Volume~II receives a state bundle
\begin{equation}
 \mathfrak S_h^{(r)}=
 (|\Psi_h\rangle_r,\cM_r^2,
 \mathfrak T_r,\mathfrak R_r,
 \{\cO_{A,r}^{\rm local}\},
 \mathsf{ConvergenceManifest}_r).
\label{eq:state-bundle}
\end{equation}
It must expose sector amplitudes without erasing:
\begin{itemize}
\item parton species and flavor;
\item color indices and coupler identity;
\item parton helicities and target helicity;
\item intrinsic momentum and angular basis labels;
\item Fock-sector identity;
\item regulator, basis, and renormalization metadata;
\item common microscopic member identity for uncertainty propagation.
\end{itemize}

\subsection{Minimum admissibility conditions}

A state may be used for central common-parent GTMD predictions only when:
\begin{enumerate}
\item it is normalized and its exact charge, color, support, plus-momentum,
and \(J^z\) ledgers close;
\item the mass operator is self-adjoint and the eigen residual is below the
declared tolerance;
\item all sectors needed for the claimed species and interference structures
are present or matched through explicit induced operators;
\item regulator, basis, Fock, and numerical convergence have been studied for
at least the calibration observables and one nontrivial withheld observable;
\item local currents/operators are matched consistently and satisfy the
applicable Ward and rotational diagnostics;
\item the state and all operator identities are fully typed;
\item no present fitted named-TMD curve is needed to define the central state;
\item the provenance graph can trace every parameter to an interaction,
current, matching term, boundary condition, or discrepancy operator.
\end{enumerate}

\begin{requirement}[Claim-scoped admissibility]
\label{req:claim-scoped}
Admissibility is claim dependent.  A \(|qqq\rangle\) pilot may be admissible
for a valence T-even benchmark while being explicitly inadmissible for central
antiquark, dynamical-gluon, or naive-\(T\)-odd predictions.  Software must
carry the allowed claim set and refuse stronger labels.
\end{requirement}

\subsection{First microscopic pilot deliverable}

The first serious nucleon pilot should solve at multiple truncations a state
containing at least the sectors in \cref{eq:minimal-nucleon-fock}, then
report:
\begin{enumerate}
\item spectrum, normalization, sector and OAM diagnostics;
\item charge, axial, magnetic, and selected energy-momentum matrix elements;
\item flavor-resolved longitudinal momentum densities;
\item quark and gluon helicity diagnostics;
\item regulator, basis, and Fock convergence;
\item the complete state bundle of \cref{eq:state-bundle} for nonzero-transfer
overlap construction.
\end{enumerate}
No GTMD-specific coefficient should be introduced in Volume~I to improve a
future projection.

\section{Acceptance criteria and immediate research decisions}
\label{sec:acceptance}

\subsection{Volume I acceptance criteria}

This volume is formally complete for implementation when:
\begin{enumerate}
\item the conventions and typed variable distinctions in
\cref{sec:front-form} are adopted without implicit aliases;
\item a complete \texttt{RegulatorSpec}, \texttt{BasisSpec}, and
\texttt{SectorId} can identify every state-space block;
\item gauge, global-color, zero-mode, and boundary data are distinct fields;
\item the mass operator is represented as typed self-adjoint blocks with
separate canonical, effective, induced, and counterterm provenance;
\item eliminating a sector automatically creates an induced-operator
obligation;
\item renormalization conditions and sector-dependent coefficients are
versioned and testable;
\item truncation levels form a directed comparison system and convergence is
reported for matched observables;
\item states carry normalization, phase, residual, and ledger manifests;
\item currents are paired with the Hamiltonian and Ward/angular diagnostics
are available;
\item the C1 identity spine enforces fail-closed composition;
\item the validation matrix V1.1--V1.18 is machine traceable;
\item the passage-to-Volume-II contract is satisfied for every claimed
central species and observable class.
\end{enumerate}

\subsection{Decisions fixed by this volume}

The following are no longer open implementation choices:
\begin{enumerate}
\item The project uses \(x^\pm=(x^0\pm x^3)/\sqrt2\) and
\(\cM^2=2P^+P^- -\PT^2\).
\item \(\bD\) and \(\bM\), and all other transverse kinds in
\cref{sec:typed-transverse}, are distinct types.
\item Transverse rank is inseparable from mass normalization, tensor basis,
Bessel order, and Fourier phase.
\item Gauge constraints, global color singlets, approximate isospin, and
front-form rotations have different formal statuses.
\item The central dynamical object is a self-adjoint regulated mass operator
on a graded physical Fock space.
\item Effective Hamiltonians require consistently transformed operators.
\item Convergence is claimed for matched observables, not bare coefficients
or Fock probabilities.
\item \(\Amp,\Dens,\Match,\Red,\Proc\) are distinct map classes.
\item C1 is adapter first and must preserve the C0 baseline byte for byte.
\end{enumerate}

\subsection{Decisions intentionally deferred}

The following require separate formal or numerical work:
\begin{enumerate}
\item the primary microscopic nucleon Hamiltonian realization and its
regulator;
\item the detailed counterterm basis and renormalization conditions for that
realization;
\item the treatment of zero modes beyond the minimum gauge-consistency datum;
\item the first complete \(|qqq\rangle\oplus|qqqg\rangle\oplus
|qqqq\bar q\rangle\) basis and computational scale;
\item the operator matching from that Hamiltonian regulator to the QCD GTMD
scheme;
\item the nuclear EFT Hamiltonian and quark-gluon-to-hadronic matching;
\item the dynamical Wilson-line expansion and rapidity subtraction.
\end{enumerate}
These are physics decisions, not missing software defaults.

\section{Conclusion}

The transition from the canonical spin-1 TMD boundary to a genuinely
predictive model begins with a well-defined state, not with another layer of
named-distribution parameterization.  This volume has specified that state's
regulated Hilbert space, particle statistics, gauge and color constraints,
front-form spin structure, self-adjoint block mass operator, effective-sector
logic, renormalization datum, directed truncation system, conserved ledgers,
and consistent current family.  It has also converted the Stage~0 audit's
most consequential findings into a fail-closed implementation contract for
coordinate, rank, sector, path, operator, and map identity.

The current accepted model remains untouched and scientifically useful.  Its
correlator parents become regression oracles and comparison data while the
new state-space machinery is introduced around them.  C1 can therefore be
implemented immediately without choosing the final microscopic Hamiltonian.
The next formal volume will use the state bundle defined here to derive
nonzero-transfer quark, antiquark, and gluon GTMD overlaps and to prove the
matched reduction relations to TMDs, GPDs, PDFs, form factors, currents, and
orbital moments.

\appendix

\section{Normalization convention and alternative adapters}
\label{app:normalization}

\subsection{Canonical Fock-state convention}

Let \(\vark=16\pi^3=2(2\pi)^3\).  A commonly used expansion is
\begin{equation}
 |P,\Lambda\rangle
 =\sum_{N,\beta}\int
 \left[\prod_{i=1}^{N}\frac{\dd x_i\,\dd^2\bm k_{iT}}
 {\sqrt{x_i}\,\vark}\right]
 \vark\,\delta\!\left(1-\sum_i x_i\right)
 \delta^{(2)}\!\left(\sum_i\bm k_{iT}\right)
 \psi_{N\beta}^{\Lambda}
 |N;\{p_i,\beta_i\}\rangle.
\label{eq:appendix-standard-expansion}
\end{equation}
With canonical creation-operator normalization, the wavefunction norm is
\begin{equation}
 \sum_{N,\beta}\int
 \left[\prod_{i=1}^{N}\frac{\dd x_i\,\dd^2\bm k_{iT}}
 {\vark}\right]
 \vark\delta\!\left(1-\sum_i x_i\right)
 \delta^{(2)}\!\left(\sum_i\bm k_{iT}\right)
 |\psi_{N\beta}^{\Lambda}|^2=1,
\label{eq:appendix-standard-norm}
\end{equation}
which is \cref{eq:fock-normalization}.  The \(x_i^{-1/2}\) factors are thus
part of the state expansion, not the probability measure.

\subsection{Unitary normalization adapter}

Suppose another implementation uses amplitudes \(\phi\) and measure
\(\dd\Gamma'\).  A normalization adapter is a multiplication/kernel operator
\(U\) satisfying
\begin{equation}
 \int\dd\Gamma\,|U\phi|^2
 =\int\dd\Gamma'\,|\phi|^2.
\label{eq:normalization-isometry}
\end{equation}
Every transformed Hamiltonian and operator must obey
\begin{equation}
 \cM'^2=U^{-1}\cM^2U,
 \qquad
 \cO'=U^{-1}\cO U.
\label{eq:normalization-operator-transform}
\end{equation}
Converting only the wavefunction while retaining old kernels is forbidden.

\section{Block-operator identities}
\label{app:block}

For an orthogonal direct sum \(\Hil=\bigoplus_\nu\Hil_\nu\), a vector is
\(\psi=(\psi_\nu)_\nu\) and
\begin{equation}
 (\cM^2\psi)_\nu=\sum_\mu\cM^2_{\nu\mu}\psi_\mu.
\label{eq:block-action}
\end{equation}
Then
\begin{align}
 \langle\phi|\cM^2\psi\rangle
 &=\sum_{\nu\mu}
 \langle\phi_\nu|\cM^2_{\nu\mu}\psi_\mu\rangle,
\label{eq:block-inner-1}\\
 \langle\cM^2\phi|\psi\rangle
 &=\sum_{\nu\mu}
 \langle\cM^2_{\mu\nu}\phi_\nu|\psi_\mu\rangle.
\label{eq:block-inner-2}
\end{align}
Equality for all \(\phi,\psi\) gives
\((\cM^2_{\nu\mu})^\dagger=\cM^2_{\mu\nu}\), with the usual domain
qualifications in infinite-dimensional spaces.

\section{Illustrative color-singlet couplers}
\label{app:color}

The normalized three-quark baryon color tensor is
\begin{equation}
 C^{(qqq)}_{ijk}=\frac{1}{\sqrt6}\epsilon_{ijk}.
\label{eq:qqq-singlet}
\end{equation}
For a quark-antiquark singlet,
\begin{equation}
 C^{(q\bar q)}_{i\bar j}=\frac{1}{\sqrt3}\delta_{i\bar j},
\label{eq:qqbar-singlet}
\end{equation}
and for a two-gluon singlet,
\begin{equation}
 C^{(gg)}_{ab}=\frac{1}{\sqrt8}\delta_{ab}.
\label{eq:gg-singlet}
\end{equation}
Multi-quark sectors generally have singlet multiplicities greater than one.
Their basis must record the recoupling tree or an invariant multiplicity
index; the label ``color singlet'' alone is not enough to reconstruct hidden
color interference.

The invariant tests are
\begin{equation}
 \left(\sum_i T_i^a\right)C=0,
 \qquad
 C^\dagger C=\Id_{\rm multiplicity},
\label{eq:color-coupler-tests}
\end{equation}
with each \(T_i^a\) in the appropriate representation.

\section{Stage 0 crosswalk and C1 requirement identifiers}
\label{app:c0-crosswalk}

The actual repository paths and existing partial implementations are
specified by:
\begin{itemize}
\item \texttt{docs/next\_level/stage0\_repository\_audit.md};
\item \texttt{docs/next\_level/stage0\_coverage\_matrix.json};
\item \texttt{docs/next\_level/stage0\_regression\_baseline.json};
\item \texttt{docs/next\_level/stageA\_migration\_plan.md};
\item \texttt{docs/next\_level/architecture\_decisions/};
\item \texttt{scripts/validate\_stage0\_architecture\_docs.py}.
\end{itemize}
This volume does not duplicate their file-by-file inventory because the audit
is the repository authority and can evolve through versioned commits.  The
following stable identifiers connect C1 to the formalism.

\begin{longtable}{L{0.15\textwidth}L{0.32\textwidth}L{0.43\textwidth}}
\toprule
ID & Formal source & C1 obligation\\
\midrule
\endhead
C1.COORD & \cref{sec:typed-transverse,eq:coordinate-spec} & Implement distinct coordinate kinds and explicit adapters.\\
C1.RANK & \cref{sec:rank-metadata,eq:rank-software-spec} & Store rank, mass, tensor, Bessel, and phase metadata around legacy transforms.\\
C1.SECTOR & \cref{eq:fock-label,eq:sector-id} & Implement immutable flavor-resolved sector identity and charge checks.\\
C1.PATH & \cref{eq:wilson-path-id,eq:gluon-link-id} & Implement path/link-pair and gluon \(f/d\) identity without implicit defaults.\\
C1.OPID & \cref{eq:decorated-operator-id} & Build decorated operator identity with explicit not-applicable versus unspecified fields.\\
C1.MAP & \cref{eq:map-class-enum,eq:composition-predicate} & Separate map classes and fail closed without a registered adapter.\\
C1.ADAPT & \cref{prin:adapter-first} & Wrap accepted objects; do not mutate central data or formulas.\\
C1.INJECT & \cref{sec:c1-spine} & Add deliberate coordinate, rank, mass, path, color, sector, scheme, and map mismatch tests.\\
C1.REGRESS & \cref{req:regression-immutability} & Pass at least the 484-test baseline and preserve all eight authoritative hashes byte for byte.\\
C1.DOC & \cref{sec:software} & Produce schemas, decision records, API documentation, and a migration manifest keyed to these IDs.\\
\bottomrule
\end{longtable}

\section{Machine-readable acceptance sketch}
\label{app:machine-readable}

A machine-readable companion may represent a requirement as
\begin{verbatim}
{
  "id": "C1.RANK",
  "formal_source": "VolumeI:req:rank-complete",
  "status": "normative",
  "objects": ["RankSpec"],
  "positive_tests": ["rank_metadata_roundtrip"],
  "negative_tests": [
    "reject_j0_for_rank1",
    "reject_missing_mass_reference",
    "reject_fourier_phase_mismatch"
  ],
  "regression": {
    "authoritative_hashes_unchanged": true
  }
}
\end{verbatim}
The JSON is an implementation companion.  The equations and definitions in
this volume remain the normative scientific source.

\small
\begin{thebibliography}{99}

\bibitem{Dirac1949}
P.~A.~M.~Dirac,
``Forms of relativistic dynamics,''
Rev. Mod. Phys. \textbf{21}, 392 (1949).

\bibitem{BrodskyPauliPinsky1998}
S.~J.~Brodsky, H.-C.~Pauli, and S.~S.~Pinsky,
``Quantum chromodynamics and other field theories on the light cone,''
Phys. Rept. \textbf{301}, 299 (1998),
arXiv:hep-ph/9705477.

\bibitem{Carbonell1998}
J.~Carbonell, B.~Desplanques, V.~A.~Karmanov, and J.-F.~Mathiot,
``Explicitly covariant light-front dynamics and relativistic few-body
systems,''
Phys. Rept. \textbf{300}, 215 (1998),
arXiv:nucl-th/9804029.

\bibitem{Miller2000}
G.~A.~Miller,
``Light front quantization: A technique for relativistic and realistic
nuclear physics,''
Prog. Part. Nucl. Phys. \textbf{45}, 83 (2000),
arXiv:nucl-th/0002059.

\bibitem{Vary2009}
J.~P.~Vary et al.,
``Hamiltonian light-front field theory in a basis function approach,''
Phys. Rev. C \textbf{81}, 035205 (2010),
arXiv:0905.1411.

\bibitem{Wilson1994}
K.~G.~Wilson, T.~S.~Walhout, A.~Harindranath, W.-M.~Zhang,
R.~J.~Perry, and S.~D.~Glazek,
``Nonperturbative QCD: A weak-coupling treatment on the light front,''
Phys. Rev. D \textbf{49}, 6720 (1994),
arXiv:hep-th/9401153.

\bibitem{GlazekWilson1993}
S.~D.~Glazek and K.~G.~Wilson,
``Renormalization of Hamiltonians,''
Phys. Rev. D \textbf{48}, 5863 (1993).

\bibitem{GlazekWilson1994}
S.~D.~Glazek and K.~G.~Wilson,
``Perturbative renormalization group for Hamiltonians,''
Phys. Rev. D \textbf{49}, 4214 (1994).

\bibitem{BrodskyHillerMcCartor1998}
S.~J.~Brodsky, J.~R.~Hiller, and G.~McCartor,
``Pauli--Villars as a nonperturbative ultraviolet regulator in discretized
light-cone quantization,''
Phys. Rev. D \textbf{58}, 025005 (1998),
arXiv:hep-th/9802120.

\bibitem{Karmanov2008}
V.~A.~Karmanov, J.-F.~Mathiot, and A.~V.~Smirnov,
``Systematic renormalization scheme in light-front dynamics with Fock space
truncation,''
Phys. Rev. D \textbf{77}, 085028 (2008),
arXiv:0801.4507.

\bibitem{Karmanov2012}
V.~A.~Karmanov, J.-F.~Mathiot, and A.~V.~Smirnov,
``Ab initio nonperturbative calculation of physical observables in
light-front dynamics: Application to the Yukawa model,''
Phys. Rev. D \textbf{86}, 085006 (2012),
arXiv:1204.3257.

\bibitem{Feshbach1958}
H.~Feshbach,
``Unified theory of nuclear reactions,''
Ann. Phys. \textbf{5}, 357 (1958).

\bibitem{Feshbach1962}
H.~Feshbach,
``A unified theory of nuclear reactions. II,''
Ann. Phys. \textbf{19}, 287 (1962).

\bibitem{Okubo1954}
S.~Okubo,
``Diagonalization of Hamiltonian and Tamm--Dancoff equation,''
Prog. Theor. Phys. \textbf{12}, 603 (1954).

\bibitem{Kato1995}
T.~Kato,
\emph{Perturbation Theory for Linear Operators},
Springer, Berlin (1995 reprint).

\bibitem{CookeMiller2002}
J.~R.~Cooke and G.~A.~Miller,
``Deuteron binding energies and form factors from light-front field theory,''
Phys. Rev. C \textbf{66}, 034002 (2002),
arXiv:nucl-th/0112037.

\bibitem{Xu2025}
S.~Xu et al.,
``Towards a first principles light-front Hamiltonian for the nucleon,''
Phys. Lett. B \textbf{867}, 139599 (2025),
arXiv:2408.11298.

\bibitem{Mondal2025}
C.~Mondal, S.~Kaur, J.~Wu, S.~Xu, X.~Zhao, and J.~P.~Vary,
``Basis light-front quantization approach to deuteron,''
arXiv:2505.12889 (2025).

\bibitem{Cao2026}
X.~Cao, S.~Cheng, Y.~Duan, Y.~Li, S.~Xu, and X.~Zhao,
``Non-perturbative flavor asymmetry in the nucleon and deuteron: The
light-front Hamiltonian effective field theory approach,''
arXiv:2601.13567 (2026).

\bibitem{Collins2011}
J.~Collins,
\emph{Foundations of Perturbative QCD},
Cambridge University Press (2011).

\bibitem{Meissner2009}
S.~Mei{\ss}ner, A.~Metz, and M.~Schlegel,
``Generalized parton correlation functions for a spin-1/2 hadron,''
JHEP \textbf{08}, 056 (2009),
arXiv:0906.5323.

\bibitem{LorcePasquini2011}
C.~Lorc\'e and B.~Pasquini,
``Quark Wigner distributions and orbital angular momentum,''
Phys. Rev. D \textbf{84}, 014015 (2011),
arXiv:1106.0139.

\bibitem{Boer2016}
D.~Boer, S.~Cotogno, T.~van Daal, P.~J.~Mulders, A.~Signori, and
Y.-J.~Zhou,
``Gluon and Wilson loop TMDs for hadrons of spin \(\le1\),''
JHEP \textbf{10}, 013 (2016),
arXiv:1607.01654.

\bibitem{Cotogno2017}
S.~Cotogno, T.~van Daal, and P.~J.~Mulders,
``Positivity bounds on gluon TMDs for hadrons of spin \(\le1\),''
JHEP \textbf{11}, 185 (2017),
arXiv:1709.07827.

\end{thebibliography}

\end{document}
